\section{Symmetry classification}

\subsection{Local characterization of symmetries}

\begin{lemma}[Common transpose sign of the conjugation factors]
    \label{thm:mpu_conjugation_common_transpose_sign}
    \lean{Matrix.transpose_mul_map_star_of_mem_unitaryGroup,
        Matrix.eq_smul_transpose_of_map_star_mul_eq_smul_one,
        Matrix.exists_common_transpose_sign_of_kronecker_map_star_mul_eq_one,
        Matrix.transpose_eq_self_or_transpose_eq_neg_of_kronecker_map_star_mul_eq_one}
    \leanok
    Let $x$ and $y$ be unitary matrices over $\mathbb C$, indexed by nonempty
    finite sets, whose entrywise conjugates satisfy
    \begin{align}
        \overline xx\otimes\overline yy & =\Id.
    \end{align}
    Then there is a sign $\sigma\in\{1,-1\}$ with
    \begin{align}
        x & =\sigma x^T, & y & =\sigma y^T.
    \end{align}
    Equivalently, either $x=x^T$ and $y=y^T$, or $x=-x^T$ and $y=-y^T$.
    In \cite[lines~1033--1038]{Cirac2017MPU}, this dichotomy is the algebraic
    conclusion.
    % Source lines: paper_v2.tex 1033--1038.
\end{lemma}

\begin{proof}\leanok
    \uses{thm:mpug_kronecker_identity_factors, thm:mpug_unitary_entrywise_self_sign}
    A Kronecker product equal to the identity has reciprocal scalar factors, so
    $\overline xx=e^{i\phi}\Id$ and $\overline yy=e^{-i\phi}\Id$. Since $x$ is
    unitary, $x^T\overline x=\Id$; multiplying $\overline xx=e^{i\phi}\Id$ on the left by $x^T$ gives $x=e^{i\phi}x^T$, and similarly
    $y=e^{-i\phi}y^T$. Transposing $x=e^{i\phi}x^T$ gives $x^T=e^{i\phi}x$, hence
    $x=e^{2i\phi}x$; since $x$ is invertible, $e^{2i\phi}=1$, so $e^{i\phi}=\pm1$ and
    $e^{-i\phi}=e^{i\phi}$. The two matrices therefore carry one transpose sign.
\end{proof}

\begin{lemma}[Conjugation symmetry]
    \label{prop1conj}
    \notready
    \discussion{5713}
    \uses{FundamentalMPU, SF}
    An MPU in standard form is invariant under complex conjugation if and only
    if there are unitaries $x$ and $y$ such that either
    $x=x^T$ and $y=y^T$, or $x=-x^T$ and $y=-y^T$, and
    \begin{align}
        \overline u & =(x\otimes y)u,                 &
        \overline v & =v(y^\dagger\otimes x^\dagger).
    \end{align}
    In \cite[Proposition~V.1]{Cirac2017MPU}, this equivalence is the unrestricted conjugation
    characterization.
    % Source lines: paper_v2.tex 1008--1039.
    % The forward result below concerns supplied reduced canonical data.
    % The contraction in lem:mpu_local_conjugation_contract is on positive
    % blocked rings; no unblocked all-length converse has been established.
\end{lemma}

\begin{lemma}[Conjugation of supplied source decompositions]
    \label{lem:mpu_source_conjugation_transport}
    \lean{MPOTensor.sourceCutM₁_mapStar, MPOTensor.sourceCutM₂_mapStar, MPOTensor.source_ranks_mapStar, MPOTensor.normalizedFlattening_mapStar, MPOTensor.IsMPUCanonicalFormII.mapStar, MPOTensor.SourceCutSVD.mapStar, MPOTensor.SourceFactors.mapStar, MPOTensor.SourceFactors.sourceU_mapStar, MPOTensor.SourceFactors.sourceV_mapStar}
    \leanok
    \uses{def:mpu_canonical_form_ii, def:mpu_weighted_source_factors,
        def:mpu_source_svd_data}
    Both source cuts satisfy
    $M_i(\overline{\mathcal U})=\overline{M_i(\mathcal U)}$, and their ranks
    are unchanged. A supplied compact decomposition
    $M_i=V_i^\dagger D_i G_i$ conjugates to
    $\overline{M_i}=\overline{V_i}^\dagger D_i\overline{G_i}$, retaining
    the positive diagonal matrices $D_i,D_i^{-1}$.
    Supplied source factors $X_i,Y_i,Z_i$ at weight $\rho$ conjugate to
    source factors of $\overline{\mathcal U}$ at weight $\overline\rho$,
    in the canonical identifications of the equal rank dimensions. Their
    gates are exactly $\overline u,\overline v$ in those coordinates.
    Normalized flattening commutes with conjugation. Reduced canonical form II
    is preserved, with the same real diagonal fixed matrix.
    These are the source-data identifications in
    \cite[Proposition~V.1, lines~1030--1033]{Cirac2017MPU}.
\end{lemma}

\begin{proof}
    \leanok
    \uses{thm:mpu_physical_adjoint_cfii, thm:mpu_physical_adjoint_full_support}
    Conjugate both factorizations, both normalization identities and both
    right-inverse identities. Products, sums, adjoints and traces commute
    with conjugation. The singular-value matrices and the diagonal positive
    fixed matrix are real. Conjugating the gate contractions therefore gives
    exactly the conjugate gates. The canonical blocks and ambient coisometry
    conjugate coefficient by coefficient, retaining dimensions and full support.
\end{proof}

\begin{lemma}[Forward conjugation equations in a reduced canonical presentation]
    \label{lem:mpu_source_local_conjugation_forward}
    \lean{MPOTensor.supplied_source_factor_gate_relations, MPOTensor.IsMPUCanonicalFormII.exists_supplied_source_gate_unitary_gauges_of_mpo_eq, MPOTensor.IsMPUCanonicalFormII.exists_local_conjugation_gauges, MPOTensor.IsMPUCanonicalFormII.exists_local_conjugation_gauges_common_transpose_case}
    \leanok
    \uses{def:mpu_canonical_form_ii, def:mpu_weighted_source_factors}
    Let $\mathcal U$ be in reduced canonical form II, with supplied source
    factors whose unshifted gate $u$ is unitary between its physical-pair and
    source-rank coordinate spaces. If
    $U^{(N)}=\overline{U^{(N)}}$ at every length $N>1$, there are unitaries
    $x,y$ such that either both are symmetric or both are skew-symmetric, and
    \begin{align}
        \overline u & =(x\otimes y)u,                 &
        \overline v & =v(y^\dagger\otimes x^\dagger).
    \end{align}
    The gates $u,v$ are the supplied gates. This is the forward argument of
    \cite[Proposition~V.1, lines~1030--1038]{Cirac2017MPU} for a reduced
    canonical presentation, with the full-support restriction recorded in
    \cite{gap:mpu_canonical_form_full_support}.
\end{lemma}

\begin{proof}
    \leanok
    \uses{lem:mpu_source_conjugation_transport, FundamentalMPU,
        lem:mpu_local_conjugation_sign, thm:mpu_source_factorizations}
    Apply the letter-level Fundamental Theorem to $\mathcal U$ and
    $\overline{\mathcal U}$. Compare the actual conjugated source factors
    with those transported by the resulting bond unitary. Normalized factor
    comparison gives two unitary changes $W_1,W_2$ of source-rank
    coordinates. With the bond unitary $z$, the source factors $X_i',Y_i'$
    of $\overline{\mathcal U}$ satisfy
    \begin{align}
        X_1' & =(\Id\otimes\overline z)X_1W_1,         &
        Y_1' & =W_1^\dagger Y_1(z^T\otimes\Id),          \\
        X_2' & =(z\otimes\Id)X_2W_2,                   &
        Y_2' & =W_2^\dagger Y_2(\Id\otimes z^\dagger).
    \end{align}
    The gate $u$ is contracted from $Y_1,Y_2$ and the gate $v$ from
    $X_1,X_2$. In both contractions $z^T$ meets $\overline z$ and $z$ meets
    $z^\dagger$ on the internal bond, so
    \begin{align}
        \overline u & =(W_2^\dagger\otimes W_1^\dagger)u, &
        \overline v & =v(W_1\otimes W_2).
    \end{align}
    Thus $x=W_2^\dagger$ and $y=W_1^\dagger$, after transport back through
    the canonical rank identifications. The unshifted equation gives the
    common transpose sign.
\end{proof}

\begin{lemma}[Transpose signs from the unshifted conjugation equation]
    \label{lem:mpu_local_conjugation_sign}
    \lean{Matrix.kronecker_map_star_mul_eq_one_of_gate_conjugation, Matrix.transpose_eq_self_or_transpose_eq_neg_of_gate_conjugation}
    \leanok
    Let $L,R,P$ be finite coordinate sets with $L,R$ nonempty, let
    $u:\mathbb C^P\to\mathbb C^L\otimes\mathbb C^R$ satisfy
    $uu^\dagger=\Id$, and let $x,y$ be unitaries on $L,R$. If
    $\overline u=(x\otimes y)u$, then
    \begin{align}
        (\overline xx)\otimes(\overline yy) & =\Id.
    \end{align}
    Consequently either $x^T=x$ and $y^T=y$, or $x^T=-x$ and $y^T=-y$.
    This is the step in \cite[lines~1033--1038]{Cirac2017MPU} starting from
    the local gate equation.
\end{lemma}

\begin{proof}
    \leanok
    \uses{thm:mpu_conjugation_common_transpose_sign}
    Conjugating the gate equation gives
    $u=(\overline x\otimes\overline y)(x\otimes y)u$.
    Multiplication by $u^\dagger$ cancels $uu^\dagger=\Id$ and yields the
    Kronecker identity. Its two factors have one common transpose sign.
\end{proof}

\begin{lemma}[Contraction of the local conjugation equations]
    \label{lem:mpu_local_conjugation_contract}
    \lean{Matrix.IsUnitaryBetween.map_star, MPOTensor.mapStar, MPOTensor.mpo_map_star, MPOTensor.TwoSiteStandardFormData.mapStar, MPOTensor.twoSiteUnshiftedLayer_map_star, MPOTensor.twoSiteShiftedLayer_map_star, MPOTensor.TwoSiteStandardFormData.map_star_mpo_eq_of_gate_conjugation}
    \leanok
    \uses{def:mpu_supplied_two_site_standard_form}
    Let $W$ have a supplied two-site standard form with unitary gates $u,v$,
    and let $x,y$ be unitaries satisfying
    \begin{align}
        \overline u & =(x\otimes y)u,                 &
        \overline v & =v(y^\dagger\otimes x^\dagger).
    \end{align}
    Conjugation of all coefficients gives a standard form of $\overline W$
    with gates $\overline u,\overline v$, and
    $\overline{W^{(N)}}=W^{(N)}$ at every positive blocked length $N$.
    This is the contraction direction of
    \cite[Proposition~V.1, lines~1029--1030]{Cirac2017MPU}, on blocked rings.
\end{lemma}

\begin{proof}
    \leanok
    \uses{thm:mpu_supplied_two_site_periodic_circuit,
        def:mpu_two_site_circuit_layers}
    Conjugation preserves the two open contractions and their unitary
    identities. It also conjugates every factor of each circuit layer.
    Let $\Gamma_N$ act on the intermediate labels by $x$ on every first
    component $h_{s,1}$ and by $y$ on every second component $h_{s,2}$.
    The two local equations give $C_{\overline u,N}=\Gamma_NC_{u,N}$ and,
    regrouping the same legs along the shifted pairs,
    $C_{\overline v,N}=C_{v,N}\Gamma_N^\dagger$. Hence
    \begin{align}
        \overline{U^{(N)}(W)}
         & =C_{\overline v,N}C_{\overline u,N}
        =C_{v,N}\Gamma_N^\dagger\Gamma_NC_{u,N}
        =C_{v,N}C_{u,N}
        =U^{(N)}(W).
        \label{eq:mpu_local_conjugation_contract_circuit}
    \end{align}
\end{proof}

\begin{lemma}[Admissible conjugation symmetry]
    \label{prop:mpu_admissible_conjugation_symmetry}
    \notready
    \discussion{5713}
    % **Scope restriction (conjugate comparison tensor):** both standard-form
    % tensors compared by Theorem~\ref{FundamentalMPU} carry
    % reduced full-support
    % source data. Documented in
    % docs/paper-gaps/mpu_canonical_form_full_support.tex.
    \uses{FundamentalMPU, SF, thm:mpu_canonical_form_ii_stabilization_converse, def:mpu_reduced_full_support_source_datum}
    Let a tensor and its conjugate comparison tensor both be equipped with
    reduced full-support source data and written in the standard forms
    determined by those data.
    The generated
    MPU is invariant under complex conjugation if and only if there are
    unitaries $x$ and $y$ such that either
    $x=x^T$ and $y=y^T$, or $x=-x^T$ and $y=-y^T$, and
    \begin{align}
        \overline u & =(x\otimes y)u,                 &
        \overline v & =v(y^\dagger\otimes x^\dagger).
    \end{align}
    For \cite[Proposition~V.1]{Cirac2017MPU}, this equivalence is the reduced-source version of the conjugation
    characterization.
    % Source lines: paper_v2.tex 1008--1039.
\end{lemma}

\begin{lemma}[Symmetric branch of the unitary congruence normal form]
    \label{lem:mpu_symmetric_unitary_congruence}
    \lean{Matrix.exists_symmetric_unitary_squareRoot}
    \leanok
    Let $x$ be a symmetric unitary matrix. There is a symmetric unitary matrix
    $S$ such that
    \begin{align}
        x & = S^T S.
    \end{align}
    In \cite[Lemma~\texttt{lemma:conjclass-normalform-continuous}]{Cirac2017MPU}, this factorization is the symmetric
    branch.
    % Source lines: paper_v2.tex 1054--1064.
\end{lemma}

\begin{proof}\leanok
    The finite spectral decomposition of $x$ takes the form
    \begin{align}
        x & = \sum_{\lambda\in\sigma(x)} \lambda P_\lambda
    \end{align}
    and set
    \begin{align}
        S & = \sum_{\lambda\in\sigma(x)} \sqrt{\lambda}\,P_\lambda.
    \end{align}
    Each spectral value has modulus one, so its square root has modulus one,
    and hence $S$ is unitary and $S^2=x$. On the finite spectrum, interpolate
    the square-root function by a polynomial $q$. Then $S=q(x)$, and
    $x^T=x$ implies $S^T=q(x^T)=q(x)=S$. Therefore $x=S^2=S^T S$.
\end{proof}

\begin{lemma}[Determinant of $i$ times a traceless Hermitian involution]
    \label{lem:mpu_det_i_hermitian_involution}
    \lean{Matrix.IsHermitian.det_I_smul_eq_one}
    \leanok
    Let $I$ be a finite set and let $H\in\operatorname{Mat}_I(\mathbb C)$ be
    Hermitian with $H^2=\Id$ and $\operatorname{tr}H=0$. Then
    \begin{align}
        \det(iH) & =1.
    \end{align}
\end{lemma}

\begin{proof}\leanok
    Let $\mu_1,\dots,\mu_{|I|}$ be the real eigenvalues of $H$. Applying
    $H^2=\Id$ to an eigenvector gives $\mu_j^2=1$, so $\mu_j=\pm1$ and
    $i\mu_j=e^{i\pi\mu_j/2}$. The trace condition gives $\sum_j\mu_j=0$.
    Hence
    \begin{align}
        \det(iH)
         & =\prod_j i\mu_j
        =\exp\Bigl(\frac{i\pi}{2}\sum_j\mu_j\Bigr)
        =1.
    \end{align}
\end{proof}

\begin{lemma}[Skew congruence from paired spectral projectors]
    \label{lem:mpu_skew_unitary_paired_projectors}
    \lean{Matrix.exists_skew_unitary_congruence_of_paired_projectors}
    \leanok
    Let $K$ and $I$ be finite sets, let $E_k\in\mathbb R$ for $k\in K$, and
    let $P_k$ be Hermitian projectors on $\mathbb C^I$. Suppose that the family
    $\{P_k,P_k^T\}_{k\in K}$ consists of mutually orthogonal projectors and
    sums to the identity. Define
    \begin{align}
        x & = \sum_{k\in K} e^{-iE_k}(P_k-P_k^T).
    \end{align}
    Then the matrices
    \begin{align}
        S & = e^{-i\pi/4}\sum_{k\in K} e^{-iE_k/2}(P_k+P_k^T),
    \end{align}
    and
    \begin{align}
        \widetilde\Lambda & = i\sum_{k\in K}(P_k-P_k^T)
    \end{align}
    are unitary, $S^T=S$, $\overline{\widetilde\Lambda}=\widetilde\Lambda$,
    $\widetilde\Lambda^T=-\widetilde\Lambda$, $\det\widetilde\Lambda=1$, and
    \begin{align}
        x & = S^T\widetilde\Lambda S.
    \end{align}
    In \cite[Lemma~\texttt{lemma:conjclass-normalform-continuous}]{Cirac2017MPU}, this is the paired-projector construction for the skew
    branch.
    % Source lines: paper_v2.tex 1066--1083.
\end{lemma}

\begin{proof}\leanok
    \uses{lem:mpu_det_i_hermitian_involution}
    Viewing the two projectors in each pair as separate members of one complete
    orthogonal family shows that a weighted sum of this family is unitary whenever every
    weight has modulus one. This applies to $S$, whose two weights at $k$
    both equal $e^{-i\pi/4}e^{-iE_k/2}$, and to
    $\widetilde\Lambda$, whose weights are $i$ and $-i$.
    Transposition exchanges the two members of each pair, proving the symmetry
    of $S$ and the skew symmetry of $\widetilde\Lambda$. Hermiticity gives
    $\overline{P_k}=P_k^T$, hence
    $\overline{\widetilde\Lambda}=\widetilde\Lambda$. Finally, orthogonality
    reduces the product to its weights, and
    \begin{align}
        \left(e^{-i\pi/4}e^{-iE_k/2}\right)^2 i & = e^{-iE_k}.
    \end{align}
    The transposed member has the negative weight, which gives
    $x=S^T\widetilde\Lambda S$. For the determinant, the same reduction gives
    $\widetilde\Lambda^2=-\Id$, so $\widetilde\Lambda=iH$ with
    $H=\sum_{k\in K}(P_k-P_k^T)$ a Hermitian involution. Since
    $\operatorname{tr}P_k^T=\operatorname{tr}P_k$, the trace of $H$ vanishes,
    and Lemma~\ref{lem:mpu_det_i_hermitian_involution} gives
    $\det\widetilde\Lambda=1$. These are the eigenvalue pairs $\pm i$ of the
    source.
\end{proof}

\begin{lemma}[Transpose pairing of spectral projectors]
    \label{lem:mpu_transpose_spectral_projector}
    \lean{Matrix.spectralProjector, Matrix.conjTranspose_spectralProjector,
        Matrix.spectralProjector_mul, Matrix.sum_spectralProjector,
        Matrix.sum_smul_spectralProjector,
        Matrix.transpose_spectralProjector_of_transpose_eq_neg}
    \leanok
    Let $x$ be a normal complex matrix with finite spectrum $\sigma(x)$, and for
    $\mu\in\mathbb C$ let $P_\mu$ be the spectral projector of $x$ at $\mu$,
    that is, the indicator function of the point $\mu$ applied to $x$. Then
    each $P_\mu$ is Hermitian,
    \begin{align}
        P_\mu P_\nu                     & =\delta_{\mu\nu}P_\mu, \\
        \sum_{\mu\in\sigma(x)}P_\mu     & =\Id,                  \\
        \sum_{\mu\in\sigma(x)}\mu P_\mu & =x.
    \end{align}
    If in addition $x^T=-x$, then
    \begin{align}
        P_\mu^T & =P_{-\mu}.
    \end{align}
    In \cite[Lemma~\texttt{lemma:conjclass-normalform-continuous}]{Cirac2017MPU}, this relation pairs $P_E$ with $P_E^T$ in the skew
    branch.
    % Source lines: paper_v2.tex 1066--1070.
\end{lemma}

\begin{proof}\leanok
    The first three identities hold for the indicator functions on
    $\sigma(x)$. For the last, interpolate the indicator of $\mu$ on the
    finite set $\sigma(x)$ by a polynomial $q$, so that $P_\mu=q(x)$. Then
    $P_\mu^T=q(x^T)=q(-x)$. Since $\sigma(x)=\sigma(x^T)=-\sigma(x)$, the
    polynomial $z\mapsto q(-z)$ agrees on $\sigma(x)$ with the indicator of
    $-\mu$, and hence $q(-x)=P_{-\mu}$.
\end{proof}

\begin{lemma}[Paired spectral data of a skew-symmetric unitary]
    \label{lem:mpu_skew_unitary_paired_spectral_data}
    \lean{Matrix.exists_paired_spectralProjectors_of_transpose_eq_neg,
        Complex.arg_nonpos_iff_not_arg_neg_nonpos}
    \leanok
    Let $x$ be a unitary matrix with $x^T=-x$. There are a finite set $K$,
    Hermitian projectors $P_k$ and numbers $E_k\in[0,\pi)$, for $k\in K$,
    such that the family $\{P_k,P_k^T\}_{k\in K}$ consists of mutually
    orthogonal projectors summing to the identity, and
    \begin{align}
        x & = \sum_{k\in K} e^{-iE_k}(P_k-P_k^T).
    \end{align}
    In \cite[Lemma~\texttt{lemma:conjclass-normalform-continuous}]{Cirac2017MPU}, this formula is the spectral decomposition in the skew
    branch.
    % Source lines: paper_v2.tex 1066--1070.
\end{lemma}

\begin{proof}\leanok
    \uses{lem:mpu_transpose_spectral_projector}
    Every spectral value of $x$ has modulus one, so it is nonzero and differs
    from its negative. By
    Lemma~\ref{lem:mpu_transpose_spectral_projector}, $\sigma(x)=-\sigma(x)$
    and $P_\lambda^T=P_{-\lambda}$. A nonzero complex number has argument in
    $(-\pi,0]$ exactly when its negative does not, so let $K$ be the set of
    spectral values with argument in $(-\pi,0]$; it contains exactly one member
    of each pair $\{\lambda,-\lambda\}$. Writing $P_\lambda$ for the projector at
    $\lambda\in K$ and $E_\lambda=-\arg\lambda\in[0,\pi)$ gives
    $e^{-iE_\lambda}=\lambda$. The members of the family
    $\{P_\lambda,P_\lambda^T\}_{\lambda\in K}$ are the spectral projectors at
    the distinct points $\pm\lambda$, so they are mutually orthogonal and sum
    to the identity. Finally,
    \begin{align}
        x & =\sum_{\lambda\in\sigma(x)}\lambda P_\lambda
        =\sum_{\lambda\in K}\lambda\left(P_\lambda-P_{-\lambda}\right)
        =\sum_{\lambda\in K}e^{-iE_\lambda}\left(P_\lambda-P_\lambda^T\right).
    \end{align}
\end{proof}

\begin{lemma}[Skew branch of the unitary congruence normal form]
    \label{lem:mpu_skew_unitary_congruence}
    \lean{Matrix.exists_skew_unitary_congruence}
    \leanok
    Let $x$ be a unitary matrix with $x^T=-x$. There are a symmetric unitary
    matrix $S$ and a real skew-symmetric unitary matrix $\widetilde\Lambda$
    with $\det\widetilde\Lambda=1$ such that
    \begin{align}
        x & = S^T\widetilde\Lambda S.
    \end{align}
    In \cite[Lemma~\texttt{lemma:conjclass-normalform-continuous}]{Cirac2017MPU}, this normal form is the skew
    branch.
    % Source lines: paper_v2.tex 1066--1083.
\end{lemma}

\begin{proof}\leanok
    \uses{lem:mpu_skew_unitary_paired_spectral_data,
        lem:mpu_skew_unitary_paired_projectors}
    Lemma~\ref{lem:mpu_skew_unitary_paired_spectral_data} writes $x$ in the
    form required by Lemma~\ref{lem:mpu_skew_unitary_paired_projectors}, which
    gives $S$ and $\widetilde\Lambda$.
\end{proof}

\begin{lemma}[Normal form of symmetric and skew-symmetric unitaries]
    \label{lemma:conjclass-normalform-continuous}
    \lean{Matrix.exists_unitary_congruence_normalForm}
    \leanok
    \discussion{5713}
    Let $x$ be unitary and satisfy $x=x^T$ or $x=-x^T$. There are a symmetric
    unitary $S$ and a matrix $\widetilde\Lambda$ such that
    \begin{align}
        x & =S^T\widetilde\Lambda S, & S & =S^T.
    \end{align}
    In the symmetric case, $\widetilde\Lambda$ is real and symmetric. In the
    skew-symmetric case, it is real and skew-symmetric. In both cases,
    $\det(\widetilde\Lambda)=1$.
    % Source lines: paper_v2.tex 1041--1084.
\end{lemma}

\begin{proof}\leanok
    \uses{lem:mpu_symmetric_unitary_congruence, lem:mpu_skew_unitary_congruence}
    In the symmetric case, Lemma~\ref{lem:mpu_symmetric_unitary_congruence}
    gives $x=S^TS$, so $\widetilde\Lambda=\Id$ works. The skew-symmetric case
    is Lemma~\ref{lem:mpu_skew_unitary_congruence}. A matrix that is both
    symmetric and skew-symmetric vanishes, and a vanishing unitary exists
    only on the zero space, where both conditions on $\widetilde\Lambda$ hold.
\end{proof}

\begin{lemma}[Odd-dimensional skew-symmetric determinant]
    \label{lem:mpu_skew_symmetric_det_zero}
    \lean{Matrix.det_eq_zero_of_transpose_eq_neg_of_odd_card}
    \leanok
    Let $I$ be a finite set of odd cardinality and let
    $A\in\operatorname{Mat}_I(\mathbb C)$ satisfy $A^T=-A$. Then
    \begin{align}
        \det A & =0.
    \end{align}
\end{lemma}

\begin{proof}\leanok
    Transpose invariance, skew symmetry, and the determinant of a negated
    matrix give
    \begin{align}
        \det A
         & =\det(A^T)
        =\det(-A)
        =(-1)^{|I|}\det A
        =-\det A.
    \end{align}
    Since the complex numbers have characteristic zero, $\det A=0$.
\end{proof}

\begin{lemma}[Parity of a nonsingular skew-symmetric matrix]
    \label{lem:mpu_skew_symmetric_even_card}
    \lean{Matrix.even_card_of_transpose_eq_neg_of_det_ne_zero}
    \leanok
    Let $I$ be finite and let $A\in\operatorname{Mat}_I(\mathbb C)$ satisfy
    $A^T=-A$ and $\det A\ne0$. Then $|I|$ is even.
\end{lemma}

\begin{proof}\leanok
    \uses{lem:mpu_skew_symmetric_det_zero}
    If $|I|$ were odd, Lemma~\ref{lem:mpu_skew_symmetric_det_zero} would give
    $\det A=0$. Hence $|I|$ is not odd and is therefore even.
\end{proof}

\begin{lemma}[Even dimension of a skew-symmetric unitary]
    \label{thm:mpu_skew_unitary_even_dimension}
    \lean{Matrix.UnitaryGroup.even_dimension_of_transpose_eq_neg}
    \leanok
    Let $U\in\operatorname{Mat}_d(\mathbb C)$ be unitary and satisfy
    $U^T=-U$. Then $d$ is even, which in \cite[Proposition~V.2, Case~II]{Cirac2017MPU} provides the parity
    obstruction.
    % Source lines: paper_v2.tex 1098--1126, especially line 1125.
\end{lemma}

\begin{proof}\leanok
    \uses{lem:mpu_skew_symmetric_even_card}
    The determinant of a unitary matrix is a unit and hence nonzero, so
    Lemma~\ref{lem:mpu_skew_symmetric_even_card} applied to $U$ shows that $d$ is even.
\end{proof}

\begin{lemma}[Invariant orthogonal planes]
    \label{lem:mpu_skew_orthogonal_plane}
    \lean{LinearIsometry.orthonormal_pair_of_sq_eq_neg,
        LinearIsometry.apply_pair_eq_standardSkew_sum,
        LinearIsometry.mapsTo_span_pair_of_sq_eq_neg,
        LinearIsometry.mapsTo_orthogonal_of_sq_eq_neg,
        LinearIsometry.restrictInvariantOrthogonal_sq,
        Matrix.orthogonalLinearIsometry_sq}
    \leanok
    \discussion{7581}
    Let $J$ be a linear isometry of a real inner-product space satisfying
    $J^2=-\Id$. For every unit vector $v$, the ordered pair $(v,-Jv)$ is
    orthonormal. Its span and its orthogonal complement are invariant under
    $J$, and the restriction to the complement also squares to $-\Id$.
    On that ordered plane the matrix of $J$ is
    $Y=\left(\begin{smallmatrix}0&1\\-1&0\end{smallmatrix}\right)$.
    Every real skew-symmetric orthogonal matrix supplies such an isometry.
\end{lemma}
\begin{proof}\leanok
    Isometry and $J^2=-\Id$ give
    $\langle x,Jy\rangle=-\langle Jx,y\rangle$. Taking $x=y=v$ shows
    $\langle v,Jv\rangle=0$. The identities $Jv=-(-Jv)$ and
    $J(-Jv)=v$ give invariance and the displayed orientation. For an invariant
    subspace $W$, the same skew-adjoint identity gives
    $\langle w,Jz\rangle=-\langle Jw,z\rangle=0$ for $w\in W$ and
    $z\in W^\perp$, proving invariance of the complement. Finally,
    $A^T=-A$ and $AA^T=\Id$ imply $A^2=-\Id$.
\end{proof}

\begin{lemma}[Sign and indexing of the standard skew form]
    \label{lem:mpu_standard_skew_indexing}
    \lean{Matrix.interleavedSkewForm_eq_neg_J_submatrix,
        Matrix.interleavedSkewFormFin_apply}
    \leanok
    \discussion{7581}
    For a finite index set $I$, identify $(i,0),(i,1)\in I\times\{0,1\}$
    with $i$ in the first and second copies of $I\sqcup I$. Under this
    identification $\Id_I\otimes Y$ corresponds to
    $\left(\begin{smallmatrix}0&\Id_I\\-\Id_I&0\end{smallmatrix}\right)$,
    the negative of the split matrix
    $\left(\begin{smallmatrix}0&-\Id_I\\\Id_I&0\end{smallmatrix}\right)$.
    For $I=\{0,\ldots,m-1\}$ the interleaved enumeration gives the same
    ordered blocks on $\{0,\ldots,2m-1\}$.
\end{lemma}
\begin{proof}\leanok
    Compare the four choices of the two binary coordinates. The two diagonal
    blocks vanish, and the off-diagonal blocks are $\Id_I$ and $-\Id_I$
    in that order. Reindexing preserves these entries.
\end{proof}

\begin{lemma}[Paired orthonormal basis]
    \label{lem:mpu_skew_paired_orthonormal_basis}
    \lean{LinearIsometry.exists_paired_orthonormalBasis}
    \leanok
    \discussion{7581}
    \uses{lem:mpu_skew_orthogonal_plane}
    Let $E$ be a real inner-product space of dimension $2m$, and let $J$ be
    a linear isometry with $J^2=-\Id$. There is an orthonormal basis
    $(e_{i,0},e_{i,1})_{i=0}^{m-1}$ such that
    $Je_{i,0}=-e_{i,1}$ and $Je_{i,1}=e_{i,0}$.
\end{lemma}
\begin{proof}\leanok
    Induct on $m$. The empty basis covers $m=0$. Otherwise choose a unit
    vector $v$ and use the ordered orthonormal pair $(v,-Jv)$.
    Its invariant orthogonal complement has dimension $2m-2$.
    Append its inductively constructed paired basis to this pair.
    The resulting orthonormal family has $2m$ members and therefore spans $E$.
\end{proof}

\begin{lemma}[Real skew-orthogonal normal form]
    \label{thm:mpu_real_skew_orthogonal_normal_form}
    \lean{Matrix.toMatrix_paired_orthonormalBasis,
        Matrix.exists_orthogonal_congruence_interleavedSkewForm}
    \leanok
    \discussion{7581}
    \uses{lem:mpu_skew_paired_orthonormal_basis, lem:mpu_standard_skew_indexing}
    Let $A\in\operatorname{Mat}_{2m}(\mathbb R)$ be orthogonal with
    $A^T=-A$. Set $\Sigma_m=\Id_m\otimes Y$, in interleaved order, where
    $Y=\left(\begin{smallmatrix}0&1\\-1&0\end{smallmatrix}\right)$.
    There is a real orthogonal matrix $O$ such that
    $O^T\Sigma_m O=A$.
    This is the real normal form used in the proof of
    \cite[Proposition~V.2]{Cirac2017MPU}.
    % Source lines: paper_v2.tex 1129--1140.
\end{lemma}
\begin{proof}\leanok
    The action of $A$ is an isometry squaring to $-\Id$. In the paired
    orthonormal basis of Lemma~\ref{lem:mpu_skew_paired_orthonormal_basis},
    its matrix is $\Sigma_m$. The orthogonal change of coordinates from
    the original basis to this paired basis gives the displayed congruence.
\end{proof}

\begin{lemma}[Complex skew-unitary normal form]
    \label{thm:mpu_complex_skew_unitary_normal_form}
    \lean{Matrix.exists_unitary_congruence_interleavedSkewForm}
    \leanok
    \discussion{7581}
    \uses{lem:mpu_skew_unitary_congruence, thm:mpu_real_skew_orthogonal_normal_form}
    Let $x\in\operatorname{Mat}_{2m}(\mathbb C)$ be unitary with $x^T=-x$.
    There is a unitary matrix $U$ such that $x=U^T\Sigma_m U$.
    This is the complex normal form used in the proof of
    \cite[Proposition~V.2]{Cirac2017MPU}.
    % Source lines: paper_v2.tex 1129--1140.
\end{lemma}
\begin{proof}\leanok
    The skew congruence factorization gives $x=S^T\Lambda S$ with $S$
    unitary and $\Lambda$ real, skew-symmetric and unitary.
    The real normal form gives $\Lambda=O^T\Sigma_m O$ with $O$ orthogonal.
    Regard $O$ as a complex unitary matrix and set $U=OS$.
\end{proof}

\begin{lemma}[Orthogonal and symplectic standard forms]
    \label{corconj}
    \notready
    \discussion{5713}
    \uses{prop1conj, lemma:conjclass-normalform-continuous, FundamentalMPU, thm:mpu_skew_unitary_even_dimension, thm:mpu_complex_skew_unitary_normal_form}
    Let an MPU in standard form be invariant under complex conjugation. Its
    standard-form unitaries can be chosen in exactly one of the following two
    cases. In Case I,
    \begin{align}
        \overline{u'} & =u', & \overline{v'} & =v',
    \end{align}
    and the choice is unique up to local orthogonal transformations. In Case II,
    \begin{align}
        \overline{u'} & =(\Sigma_\ell\otimes\Sigma_r)u', &
        \overline{v'} & =v'(\Sigma_r\otimes\Sigma_\ell),
    \end{align}
    where
    \begin{align}
        \Sigma_x & =\Id_{x/2}\otimes Y,                    &
        Y        & =\begin{pmatrix}0&1\\-1&0\end{pmatrix}.
    \end{align}
    The choice is unique up to local symplectic transformations $S_x$
    satisfying $S_x^T\Sigma_xS_x=\Sigma_x$. Case II can occur only when $r$ and
    $\ell$ are even.
    % Source lines: paper_v2.tex 1094--1163.
\end{lemma}

\begin{lemma}[Admissible orthogonal and symplectic standard forms]
    \label{prop:mpu_admissible_conjugation_standard_forms}
    \notready
    \discussion{5713}
    % **Scope restriction (standard-form comparisons):** every alternative
    % standard-form tensor compared by the fundamental theorem carries its own
    % reduced full-support source datum. Documented in
    % docs/paper-gaps/mpu_canonical_form_full_support.tex.
    \uses{prop:mpu_admissible_conjugation_symmetry, lemma:conjclass-normalform-continuous, FundamentalMPU, thm:mpu_canonical_form_ii_stabilization_converse, def:mpu_reduced_full_support_source_datum}
    Let an MPU in the standard form determined by its given source datum
    be invariant under complex
    conjugation, and
    suppose every standard-form tensor compared below carries reduced
    full-support source data. Its
    standard-form unitaries can be chosen in exactly one of the following two
    cases. In Case I,
    \begin{align}
        \overline{u'} & =u', & \overline{v'} & =v',
    \end{align}
    and the choice is unique up to local orthogonal transformations. In Case II,
    \begin{align}
        \overline{u'} & =(\Sigma_\ell\otimes\Sigma_r)u', &
        \overline{v'} & =v'(\Sigma_r\otimes\Sigma_\ell),
    \end{align}
    where
    \begin{align}
        \Sigma_x & =\Id_{x/2}\otimes Y,                    &
        Y        & =\begin{pmatrix}0&1\\-1&0\end{pmatrix}.
    \end{align}
    The choice is unique up to local transformations $S_x$ satisfying
    $S_x^T\Sigma_xS_x=\Sigma_x$. Case II can occur only when $r$ and $\ell$
    are even.
    % Source lines: paper_v2.tex 1094--1163.
\end{lemma}

\begin{lemma}[The four-dimensional symplectic congruence]
    \label{lem:mpu_symplectic_four_congruence}
    \lean{Matrix.symplecticY, Matrix.symplecticOrthogonalR,
        Matrix.symplecticOrthogonalR_identities}
    \leanok
    Let $Y=\begin{pmatrix}0&1\\-1&0\end{pmatrix}$ and
    $R=\frac{1-i}{2}[\Id_4+i(Y\otimes Y)]$, with the two tensor factors in
    that order. Then
    \begin{align}
        R^\dagger R      & =\Id_4, \\
        R^T(Y\otimes Y)R & =\Id_4.
    \end{align}
    % Source: arXiv:1703.09188, paper_v2.tex lines 1165--1185.
\end{lemma}

\begin{proof}
    \leanok
    Direct multiplication of the four entries of $Y$ and the corresponding
    entries of $Y\otimes Y$ gives both identities.
\end{proof}

\begin{lemma}[Reduction of the symplectic case to orthogonal unitaries]
    \label{cor:sympl-to-orthogonal}
    \notready
    \discussion{5713}
    \uses{corconj, lem:mpu_symplectic_four_congruence}
    In Case II of Lemma~\ref{corconj}, one can write
    \begin{align}
        u' & =(\Id_{r\ell/4}\otimes R)\widetilde u, &
        v' & =\widetilde v(\Id_{r\ell/4}\otimes R),
    \end{align}
    where $\widetilde u$ and $\widetilde v$ are orthogonal and
    \begin{align}
        R & =\frac{1-i}{2}[\Id_4+i(Y\otimes Y)].
    \end{align}
    % Source lines: paper_v2.tex 1165--1185.
\end{lemma}

\begin{lemma}[Admissible reduction of the symplectic case to orthogonal unitaries]
    \label{cor:mpu_admissible_symplectic_to_orthogonal}
    \notready
    \discussion{5713}
    % **Scope restriction (source-derived standard form):** Case II refers to the
    % source-data-restricted standard form of Lemma~\ref{prop:mpu_admissible_conjugation_standard_forms}.
    % Documented in docs/paper-gaps/mpu_canonical_form_full_support.tex.
    \uses{prop:mpu_admissible_conjugation_standard_forms, def:mpu_reduced_full_support_source_datum}
    In the admissible Case II of Lemma~\ref{prop:mpu_admissible_conjugation_standard_forms}, one can write
    \begin{align}
        u' & =(\Id_{r\ell/4}\otimes R)\widetilde u, &
        v' & =\widetilde v(\Id_{r\ell/4}\otimes R),
    \end{align}
    where $\widetilde u$ and $\widetilde v$ are orthogonal and
    \begin{align}
        R & =\frac{1-i}{2}[\Id_4+i(Y\otimes Y)].
    \end{align}
    % Source lines: paper_v2.tex 1165--1185.
\end{lemma}

\begin{lemma}[Local characterization of time reversal]
    \label{PropChiral}
    \notready
    \discussion{5714}
    \uses{FundamentalMPU, SF, def:index}
    Let an MPU be in standard form, with one-site physical space
    $\mathbb C^{d_0}\otimes\mathbb C^{d_0}$. Time-reversal symmetry forces
    $\operatorname{ind}(\mathcal U)=0$. The MPU satisfies
    $U^{(N)}=U^{(N)\dagger}$ for every even $N$ if and only if there is a
    unitary $x$ such that
    \begin{align}
        (\Id_1\otimes v_{23}\otimes\Id_4)(u_{12}\otimes u_{34})
         & =\pm(x_{12}\otimes x_{34}^\dagger)
        (\Id_1\otimes v_{23}^\dagger\otimes\Id_4).
    \end{align}
    % Source lines: paper_v2.tex 1201--1252.
\end{lemma}

\begin{proof}
    \uses{FundamentalMPU, index-well-defined, def:index,
        thm:mpu_physical_adjoint_source_ranks}
    Under time reversal, the two-site block $\mathcal U_2$ and its
    physical-adjoint comparison tensor $\mathcal U_2^\sharp$ generate the
    same MPU. The fundamental theorem identifies their source spaces,
    so their corresponding source-cut ranks agree. Physical adjunction exchanges
    the two cuts by Lemma~\ref{thm:mpu_physical_adjoint_source_ranks}. Hence
    \begin{align}
        r(\mathcal U_2)
         & =r(\mathcal U_2^\sharp)
        =\ell(\mathcal U_2).
    \end{align}
    Definition~\ref{def:index} therefore gives
    $\operatorname{ind}(\mathcal U_2)=0$, and blocking independence gives
    $\operatorname{ind}(\mathcal U)=0$.

    Applying the two gauge equations of the fundamental theorem to the standard
    forms of $\mathcal U_2$ and $\mathcal U_2^\sharp$ and substituting the second
    equation into the first gives
    \begin{align}
        x_{12}y_{12}\otimes y_{34}x_{34} & =\Id.
    \end{align}
    Therefore $y=\pm x^\dagger$, which yields the displayed local relation.
    Conversely, contract that relation and its adjoint alternately around an
    even periodic chain.
    % Source lines: paper_v2.tex 1201--1252.
\end{proof}

\begin{lemma}[Source cuts of a supplied square standard form]
    \label{lem:mpu_supplied_standard_form_source_cuts}
    \lean{MPOTensor.TwoSiteStandardFormData.cutX₁, MPOTensor.TwoSiteStandardFormData.cutY₁, MPOTensor.TwoSiteStandardFormData.cutX₂, MPOTensor.TwoSiteStandardFormData.cutY₂, MPOTensor.TwoSiteStandardFormData.sourceCutM₁_eq_cutX₁_mul_cutY₁, MPOTensor.TwoSiteStandardFormData.sourceCutM₂_eq_cutX₂_mul_cutY₂, MPOTensor.TwoSiteStandardFormData.cutY₂_cutY₁_eq_staircaseMiddle_mul_kronecker, MPOTensor.TwoSiteStandardFormData.cutX₁_cutX₂_eq_staircaseMiddle, MPOTensor.IsMPUCanonicalFormII.twoSite_cutX₁_weighted_isometry, MPOTensor.IsMPUCanonicalFormII.twoSite_cutX₂_isIsometry, MPOTensor.IsMPUCanonicalFormII.twoSite_sourceRanks, MPOTensor.IsMPUCanonicalFormII.twoSiteSourceFactors, MPOTensor.IsMPUCanonicalFormII.sourceRanks_eq_phys_of_standardForm_even_mpo_eq_adjoint, MPOTensor.IsMPUCanonicalFormII.blockTwo, MPOTensor.IsMPUCanonicalFormII.isMPUSimple_blockTwo, MPOTensor.IsMPUCanonicalFormII.twoSiteSourceFactors_sourceU, MPOTensor.IsMPUCanonicalFormII.twoSiteSourceFactors_sourceV, MPOTensor.blockTensor_one_eq_reindexPhysical}
    \leanok
    Let $W$ have a supplied two-site standard form with unitary gates
    $u:\mathbb C^d\otimes\mathbb C^d\to\mathbb C^\ell\otimes\mathbb C^r$ and
    $v:\mathbb C^r\otimes\mathbb C^\ell\to\mathbb C^d\otimes\mathbb C^d$,
    and half factors $X_1,X_2$.
    The two source cuts of $W$ itself admit factors through
    $\mathbb C^d\otimes\mathbb C^r$ and
    $\mathbb C^\ell\otimes\mathbb C^d$, respectively, given by
    \begin{align}
        \widehat X_1((i_1,i_2),\gamma;a,s)
         & =\delta_{i_1,a}X_1(i_2,\gamma;s), \\
        \widehat Y_1(a,s;\alpha,j)
         & =\sum_t X_2(\alpha,a;t)u(t,s;j),  \\
        \widehat X_2(\alpha,(i_1,i_2);t,b)
         & =X_2(\alpha,i_1;t)\delta_{i_2,b}, \\
        \widehat Y_2(t,b;j,\gamma)
         & =\sum_s u(t,s;j)X_1(b,\gamma;s).
    \end{align}
    Thus $\mathcal M_a(W)=\widehat X_a\widehat Y_a$ for $a=1,2$.
    Contracting $\widehat Y_2$ with $\widehat Y_1$ gives
    $v_{23}(u_{12}\otimes u_{34})$; contracting $\widehat X_1$ with
    $\widehat X_2$ gives $v_{23}$. These are the four-original-spin
    matrices in \cite[lines~1237--1252]{Cirac2017MPU}, without first
    blocking $W$ again. When this standard form is constructed from a simple
    canonical tensor with boundary $\rho$, its coarse left factors satisfy
    $\widehat X_1^\dagger(\Id_{d^2}\otimes\rho)\widehat X_1=\Id$ and
    $\widehat X_2^\dagger\widehat X_2=\Id$. In this canonical construction,
    the actual cut ranks are $d r$ and $\ell d$. The right factors have
    full row rank, hence right inverses; together these give normalized
    source factors at the same boundary $\rho$. If the original canonical
    tensor $U$ is simple and its standard-form tensor $W$ satisfies
    $W^{(N)}=W^{(N)\dagger}$ at every even length $N>1$, then
    $r(U)=\ell(U)=d$.
\end{lemma}

\begin{proof}\leanok
    \uses{SF}
    Substitute the open standard-form contraction in each source cut.
    Sum the two spectator Kronecker deltas. In the right-factor contraction,
    the virtual sum of $X_1$ and $X_2$ is $v$; the two remaining factors are
    $u$. The left-factor contraction uses the same virtual sum and leaves
    the first and fourth spins fixed. The first coarse left factor is
    $\Id_d\otimes X_1$ and the second is $X_2\otimes\Id_d$, after
    the stated coordinate regroupings. Their normalizations follow from
    those of the original source factors by taking these tensor products.
    Simplicity persists under blocking; the rank-product saturation argument
    gives the ranks $dr$ and $\ell d$. Each cut factors through its displayed
    pair space, so its right factor has full row rank and a right inverse.
    For the last assertion, adjoint invariance at even standard-form lengths
    identifies the even periodic families of the two-site block and its
    adjoint. Their cut ranks agree by the sign-sensitive virtual comparison, while
    adjunction exchanges the two cuts. Thus $dr(U)=\ell(U)d$, so
    $r(U)=\ell(U)$. Simplicity gives $r(U)\ell(U)=d^2$, hence both equal $d$.
\end{proof}

\begin{lemma}[Minimality of the supplied standard-form cuts]
    \label{lem:mpu_supplied_half_factor_minimality}
    \lean{MPOTensor.TwoSiteStandardFormData.X₁_mulVec_injective, MPOTensor.TwoSiteStandardFormData.X₂_mulVec_injective, MPOTensor.TwoSiteStandardFormData.exists_left_inverses, MPOTensor.TwoSiteStandardFormData.cutX₁_mulVec_injective, MPOTensor.TwoSiteStandardFormData.cutX₂_mulVec_injective, MPOTensor.TwoSiteStandardFormData.exists_cutY_right_inverses, MPOTensor.TwoSiteStandardFormData.cut_factor_ranks, MPOTensor.TwoSiteStandardFormData.sourceRanks_eq, MPOTensor.TwoSiteStandardFormData.fourSpinGate_isUnitaryBetween}
    \leanok
    Let $W$ have any supplied standard form as in
    Lemma~\ref{lem:mpu_supplied_standard_form_source_cuts}, with both gates
    unitary and positive physical, bond, and intermediate dimensions.
    The half factors $X_1,X_2$ have independent columns and left inverses.
    The displayed cut factors satisfy
    \begin{align}
        \operatorname{rank}\widehat X_1
        =\operatorname{rank}\widehat Y_1 & =dr,     \\
        \operatorname{rank}\widehat X_2
        =\operatorname{rank}\widehat Y_2 & =\ell d.
    \end{align}
    The four-spin matrix $v_{23}(u_{12}\otimes u_{34})$ is unitary.
    There are matrices $Z_1,Z_2$ with
    $\widehat Y_1Z_1=\Id_{dr}$ and $\widehat Y_2Z_2=\Id_{\ell d}$.
    Consequently the actual source ranks are $r(W)=dr$ and
    $\ell(W)=\ell d$. Neither a canonical presentation nor a normalization
    of the half factors is required. These are consequences of the supplied
    contraction in \cite[lines~532--543, 603--617 and
        1237--1244]{Cirac2017MPU}.
\end{lemma}

\begin{proof}\leanok
    \uses{lem:mpu_supplied_standard_form_source_cuts}
    The supplied contraction of the half factors is
    \begin{align}
        v(i,j;s,t) & =\sum_\beta X_1(i,\beta;s)X_2(\beta,j;t).
        \label{eq:mpu_supplied_half_factor_minimality_v}
    \end{align}
    Let $X_1z=0$, fix an intermediate coordinate $t_0$, and put
    $w(s,t)=z(s)\delta_{t,t_0}$. By
    \eqref{eq:mpu_supplied_half_factor_minimality_v},
    \begin{align}
        (vw)(i,j) & =\sum_\beta X_2(\beta,j;t_0)
        \sum_sX_1(i,\beta;s)z(s)=0.
    \end{align}
    Since $v^\dagger v=\Id$, this gives $w=v^\dagger vw=0$, hence $z=0$:
    $X_1$ is injective. Exchanging the two physical and intermediate legs
    of \eqref{eq:mpu_supplied_half_factor_minimality_v} proves the same
    assertion for $X_2$. By the formulas of
    Lemma~\ref{lem:mpu_supplied_standard_form_source_cuts},
    $\widehat X_1=\Id_d\otimes X_1$ and $\widehat X_2=X_2\otimes\Id_d$
    up to the order of the tensor factors, so both are injective.

    Write $A=v_{23}(u_{12}\otimes u_{34})$. It is unitary, and by
    Lemma~\ref{lem:mpu_supplied_standard_form_source_cuts}
    \begin{align}
        A^\dagger(j,j';(t,b),(a,s))
         & =\sum_\gamma\overline{\widehat Y_2(t,b;j,\gamma)}\,
        \overline{\widehat Y_1(a,s;\gamma,j')}.
    \end{align}
    This is \eqref{eq:mpu_supplied_half_factor_minimality_v} for the
    isometry $A^\dagger$, with half factors $\widehat Y_2^\dagger$ and
    $\widehat Y_1^\dagger$. The kernel argument above therefore makes
    $\widehat Y_1^\dagger$ and $\widehat Y_2^\dagger$ injective.
    Their left inverses, after taking adjoints, give the asserted right
    inverses. All four factor ranks follow. Multiplying each source cut
    by its right-factor inverse recovers its left factor, so the source
    cut and its left factor have equal rank.
\end{proof}

\begin{lemma}[Supplied four-spin gates and normalized source factorizations]
    \label{lem:mpu_supplied_four_spin_source_comparison}
    \lean{MPOTensor.TwoSiteStandardFormData.isMPUSimple, MPOTensor.TwoSiteStandardFormData.exists_fourSpin_source_gate_gauges, MPOTensor.TwoSiteStandardFormData.exists_fourSpin_selected_source_gate_gauges, MPOTensor.TwoSiteStandardFormData.exists_fourSpin_selected_source_gate_gauges_of_square_ranks}
    \leanok
    Let $W$ have any supplied two-site standard form with unitary gates
    $u,v$ and positive dimensions. Write
    $A=v_{23}(u_{12}\otimes u_{34})$ and $B=v_{23}$.
    Take another compact factorization of each actual source cut, normalized
    at a boundary matrix $\rho$ by
    $\widetilde X_1^\dagger(I\otimes\rho)\widetilde X_1=I$ and
    $\widetilde X_2^\dagger\widetilde X_2=I$, with right inverses of both
    right factors. Suppose its unshifted gate $\widetilde u$ is unitary.
    Then there are unitaries $x,y$ on the actual left and right source
    spaces such that, under the canonical finite-dimensional coordinate
    identifications of Lemma~\ref{lem:mpu_supplied_half_factor_minimality},
    \begin{align}
        \widetilde u & =(x\otimes y)A,                 &
        \widetilde v & =B(y^\dagger\otimes x^\dagger).
    \end{align}
    In particular, when $W$ is explicitly given in reduced canonical
    form II, it is simple and its selected normalized source factors
    satisfy these equations. No normalization of the supplied half factors
    is required. The ambient canonical presentation in this last assertion
    is an assumption, not a consequence of the displayed cut ranks.
    These are the actual four-spin contractions of
    \cite[lines~1237--1252]{Cirac2017MPU}, compared by the raw factorization
    uniqueness lemma in \cite{FrancoRubio2025MPUGauging}.
\end{lemma}

\begin{proof}\leanok
    \uses{lem:mpu_supplied_half_factor_minimality}
    The two supplied cut factorizations have full intermediate rank and
    right inverses. Compare them with the normalized factors by the raw
    decomposition uniqueness theorem. Its reciprocal positive scalars
    cancel in both contractions. The remaining two unitary factors give
    the displayed equations. If the ambient tensor is in canonical form II,
    the actual cut ranks and $\ell r=d^2$ give the simplicity criterion.
    Its selected source gates are therefore unitary and the comparison
    applies to that normalized compact factorization.
\end{proof}

\begin{lemma}[Square intermediate dimensions for any supplied adjoint-symmetric standard form]
    \label{lem:mpu_supplied_standard_form_adjoint_zero_index}
    \lean{MPOTensor.TwoSiteStandardFormData.intermediateRanks_eq_phys_of_even_mpo_eq_adjoint}
    \leanok
    Let $W$ have any supplied standard form with unitary gates on pairs
    of $d$-dimensional sites and intermediate dimensions $\ell,r>0$.
    Suppose the ambient tensor $W$ is explicitly in reduced canonical
    form II and $W^{(N)}=W^{(N)\dagger}$ for every even $N>1$.
    Then $r=\ell=d$. Only the stated even standard-form lengths are used.
    This is the zero-index comparison of
    \cite[lines~1201--1252]{Cirac2017MPU} in the explicit canonical ambient scope.
\end{lemma}

\begin{proof}\leanok
    \uses{lem:mpu_supplied_half_factor_minimality}
    The even periodic families of $W$ and its actual adjoint agree.
    Their actual source ranks agree by the sign-sensitive virtual
    comparison, while adjunction exchanges the cuts. Thus
    $\ell d=\ell(W)=r(W)=dr$. Positivity of $d$ gives $\ell=r$.
    The unitary gate $u$ also gives $\ell r=d^2$, so both intermediate
    dimensions equal $d$.
\end{proof}

\begin{lemma}[Local time reversal for any supplied standard form in a canonical ambient tensor]
    \label{lem:mpu_supplied_forward_local_time_reversal}
    \lean{MPOTensor.TwoSiteStandardFormData.exists_timeReversal_certificate_of_square_ranks, MPOTensor.TwoSiteStandardFormData.exists_timeReversal_certificate, MPOTensor.TwoSiteStandardFormData.mpo_even_eq_adjoint_of_timeReversal_certificate, MPOTensor.TwoSiteStandardFormData.even_timeReversal_iff_exists_certificate}
    \leanok
    Let $W$ be explicitly in reduced canonical form II and let $u,v$
    be any supplied standard-form gates on pairs of $d$-dimensional sites.
    Suppose $W^{(N)}=W^{(N)\dagger}$ for every even $N>1$.
    Then the intermediate dimensions satisfy $r=\ell=d$.
    Under these derived coordinate identifications there exist a unitary
    $x$ on $\mathbb C^d\otimes\mathbb C^d$ and $\sigma\in\{1,-1\}$ such that
    \begin{align}
        v_{23}(u\otimes u) & =\sigma(x\otimes x^\dagger)v_{23}^\dagger.
    \end{align}
    Conversely, these derived-rank and local-equation conditions imply
    $W^{(N)}=W^{(N)\dagger}$ at every even $N>1$.
    Thus the two conditions are equivalent. No normalization of the supplied
    half factors is assumed. The explicit ambient canonicality and
    even-length restrictions remain;
    this is the corresponding restricted equivalence of
    \cite[lines~1201--1252]{Cirac2017MPU}, not the unrestricted source criterion.
\end{lemma}

\begin{proof}\leanok
    \uses{lem:mpu_supplied_standard_form_adjoint_zero_index,
        lem:mpu_supplied_four_spin_source_comparison,
        lem:mpu_normalized_source_forward_adjoint_certificate}
    Derive $r=\ell=d$ and choose the actual normalized source factors of $W$.
    Write $A=v_{23}(u\otimes u)$ and $B=v_{23}$.
    Their normalized gates satisfy
    $\widetilde A=(X\otimes Y)A$ and
    $\widetilde B=B(Y^\dagger\otimes X^\dagger)$.
    The normalized source certificate gives
    $\widetilde A=\sigma(z\otimes z^\dagger)\widetilde B^\dagger$.
    Multiply by $X^\dagger\otimes Y^\dagger$ to obtain the displayed equation
    with $x=X^\dagger zY$. All coordinate changes are the actual source-rank
    and physical-pair identifications, and no further blocking is used.
    For the converse, the inverse comparison gives a normalized source
    certificate with unitary $z=XxY^\dagger$. Its actual even-ring
    contraction yields the asserted periodic adjoint symmetry.
\end{proof}

\begin{lemma}[Normalized physical-adjoint source factors]
    \label{lem:mpu_normalized_physical_adjoint_source_factors}
    \lean{MPOTensor.IsMPUCanonicalFormII.reindexPhysical, MPOTensor.IsMPUCanonicalFormII.sourceFactors_sourceU_isUnitaryBetween, MPOTensor.IsMPUCanonicalFormII.sourceFactors_Y₁_mul_conjTranspose, MPOTensor.IsMPUCanonicalFormII.sourceFactors_Y₂_weighted_mul_conjTranspose, MPOTensor.SourceFactors.physicalAdjoint, MPOTensor.SourceFactors.sourceU_physicalAdjoint, MPOTensor.SourceFactors.sourceV_physicalAdjoint, MPOTensor.IsMPUCanonicalFormII.twoSiteSourceFactors_physicalAdjoint}
    \leanok
    \uses{thm:mpug_canonical_adjoint_simple}
    Let $U$ be simple and in reduced canonical form II, with recorded positive
    diagonal boundary $\rho$. Bijective physical relabeling preserves this
    presentation and boundary. Physical adjunction preserves simplicity.
    For every normalized source factorization at this same boundary,
    \begin{align}
        Y_1Y_1^\dagger                   & =(d/r)\Id,    &
        Y_2(\Id_d\otimes\rho)Y_2^\dagger & =(d/\ell)\Id.
    \end{align}
    Its unshifted gate is unitary. If $r=\ell=d$, source factors for the
    physical adjoint at the same boundary are
    \begin{align}
        (X_1',Y_1',Z_1') & =(Y_2^\dagger,X_2^\dagger,X_2), \\
        (X_2',Y_2',Z_2') & =(Y_1^\dagger,X_1^\dagger,
        (\Id_d\otimes\rho)X_1),
    \end{align}
    after the canonical identifications of the exchanged cut ranks.
    Their gates are exactly $u'=v^\dagger$ and $v'=u^\dagger$ in these
    coordinates. For the actual coarse record of
    Lemma~\ref{lem:mpu_supplied_standard_form_source_cuts}, adjoint symmetry
    at even standard-form lengths derives all these prerequisites, giving an
    actual adjoint coarse record without extra normalization witnesses.
    These are the normalized cut constructions underlying
    \cite[lines~1237--1252]{Cirac2017MPU}; the Gram formulas are
    \cite[lines~1326--1398]{FrancoRubio2025MPUGauging}.
\end{lemma}

\begin{proof}\leanok
    \uses{lem:mpu_supplied_standard_form_source_cuts}
    Simplicity is the equality $r\ell=d^2$ and adjunction exchanges $r,\ell$.
    For the Gram formulas, compare a supplied cut factorization $XY$ with
    the selected normalized one. The right inverse makes their left-factor
    ranges equal. Their common weighted normalization gives a coisometry
    $A$ with $Y'=AY$, so each scalar Gram matrix is unchanged.
    At $r=\ell=d$ both Grams are identities. Taking adjoints exchanges the
    two cut factorizations. The displayed old left-factor normalizations
    give the new right inverses, and their old Gram identities give the new
    left-factor normalizations. Direct contraction gives the two adjoint gates.
\end{proof}

\begin{lemma}[Scalar transport of normalized source factors]
    \label{lem:mpu_source_factors_scalar_transport}
    \lean{MPOTensor.source_ranks_smul, MPOTensor.SourceFactors.smul, MPOTensor.SourceFactors.sourceU_smul, MPOTensor.SourceFactors.sourceV_smul}
    \leanok
    For a nonzero scalar $c$, the two source ranks of $cU$ equal those of $U$.
    Any normalized source factorization of $U$ transports to one of $cU$
    with the same weight, by keeping both left factors, multiplying both
    right factors by $c$, and multiplying both right inverses by $c^{-1}$.
    In the canonical rank identifications the transported gates are $c^2u$
    and $v$. In particular a sign disappears in the two-letter source gate.
    These are the contractions of
    \cite[lines~479--524]{Cirac2017MPU} used in the even-family comparison.
\end{lemma}

\begin{proof}\leanok
    A nonzero scalar does not change matrix rank. The two cut products acquire
    the scalar, the two left-factor normalizations stay unchanged, and the
    right-inverse products acquire $cc^{-1}=1$. The unshifted gate contracts
    two right factors, so it acquires $c^2$; the shifted gate contracts the
    unchanged left factors.
\end{proof}

\begin{lemma}[Even-family virtual rigidity up to a sign]
    \label{lem:mpu_even_family_virtual_sign_gauge}
    \lean{MPOTensor.IsMPUCanonicalFormII.exists_unitary_virtual_sign_gauge_of_mpo_eq_even}
    \leanok
    Let $U,V$ be reduced canonical-form-II MPU tensors with equal periodic
    operators at every even length greater than one. Their bond dimensions
    agree, and there are a virtual unitary $z$ and a sign $\eta\in\{1,-1\}$
    such that $V^{ij}=\eta zU^{ij}z^\dagger$ for every physical pair $i,j$.
    No equality at odd lengths is inferred. This is the virtual comparison
    needed in the even-length setting of
    \cite[lines~1209--1252]{Cirac2017MPU}; reconstructing the four-spin
    gates remains a separate step. The presentation restriction is recorded
    in \cite{gap:mpu_canonical_form_full_support}.
\end{lemma}

\begin{proof}\leanok
    \uses{thm:cpsv_normal_overlap_dichotomy}
    Equality along the even subsequence makes the cross-overlap equal the
    self-overlap there, whose limit is one. The normal-tensor overlap
    dichotomy therefore rules out vanishing cross-overlap and gives an
    unblocked invertible gauge with a unit phase $\eta$. The nonzero periodic
    vector at length two forces $\eta^2=1$. Left canonical normalization and
    irreducibility replace the invertible gauge by a unitary without changing
    the phase.
\end{proof}

\begin{lemma}[Supplied source gates from even-family equality]
    \label{lem:mpu_even_family_supplied_source_gauges}
    \lean{MPOTensor.IsMPUCanonicalFormII.eq_sign_virtualSandwich_of_mpo_eq_even, MPOTensor.IsMPUCanonicalFormII.source_ranks_eq_of_mpo_eq_even, MPOTensor.IsMPUCanonicalFormII.exists_supplied_source_gate_unitary_gauges_of_mpo_eq_even}
    \leanok
    Let $U,V$ be reduced canonical-form-II MPU tensors with the same bond
    dimension and equal periodic operators at every even length greater than
    one. For any supplied normalized source factors whose unshifted gates
    are unitary, their source ranks agree and their actual gates satisfy
    $u_V=(x\otimes y)u_U$ and $v_V=v_U(y^\dagger\otimes x^\dagger)$
    for unitary rank-coordinate transformations $x,y$, in the canonical
    identifications of equal ranks. The signed target is a virtual unitary
    conjugate of the original tensor. This is the even-family source
    comparison used in \cite[lines~1237--1252]{Cirac2017MPU}.
\end{lemma}

\begin{proof}\leanok
    \uses{lem:mpu_even_family_virtual_sign_gauge,
        lem:mpu_source_factors_scalar_transport, FundamentalMPU}
    Write $\eta V=zUz^\dagger$ with $\eta^2=1$. Transport the supplied
    factors of $V$ to those of $\eta V$. Their gates are unchanged since
    the unshifted gate acquires $\eta^2$ and the shifted gate acquires no
    scalar. Apply the normalized factor comparison under the actual virtual
    sandwich, then transport the rank-coordinate unitaries back. Nonzero
    scaling and unitary virtual conjugation preserve both source ranks.
\end{proof}

\begin{lemma}[Forward adjoint certificate for actual normalized source factors]
    \label{lem:mpu_normalized_source_forward_adjoint_certificate}
    \lean{MPOTensor.IsMPUCanonicalFormII.exists_supplied_source_adjoint_certificate, Matrix.reindex_adjoint_certificate}
    \leanok
    Let $W$ be simple and explicitly in reduced canonical form II, with
    $r(W)=\ell(W)=q$, where $q$ is its physical dimension.
    Take any actual normalized source factorization at its recorded boundary.
    If $W^{(N)}=W^{(N)\dagger}$ for every even $N>1$, then, in these rank
    identifications, its gates obey
    $A=\sigma(x\otimes x^\dagger)B^\dagger$ for a unitary $x$ and
    $\sigma\in\{1,-1\}$. The equation is preserved under every common
    coordinate relabeling of the two rank spaces and physical spaces.
    This is the normalized-factor step of
    \cite[lines~1237--1252]{Cirac2017MPU} in the stated canonical scope.
\end{lemma}

\begin{proof}\leanok
    \uses{lem:mpu_normalized_physical_adjoint_source_factors,
        lem:mpu_even_family_supplied_source_gauges}
    Construct the actual physical-adjoint factors and compare the even
    periodic families. The two resulting unitary gate equations imply
    $YX\otimes XY=I$, hence $Y=\sigma X^\dagger$ with $\sigma^2=1$.
\end{proof}

\begin{lemma}[Forward local time reversal for the constructed canonical standard form]
    \label{lem:mpu_canonical_forward_local_time_reversal}
    \lean{MPOTensor.IsMPUCanonicalFormII.exists_coarse_timeReversal_certificate, MPOTensor.IsMPUCanonicalFormII.exists_local_timeReversal_certificate, MPOTensor.IsMPUCanonicalFormII.exists_twoSiteStandardFormData_timeReversal_certificate, MPOTensor.coarse_gates_eq_four_spin_gates}
    \leanok
    Let $U$ be a simple reduced canonical-form-II tensor. Construct its actual
    source standard form $W$ with gates $u,v$, and suppose that
    $W^{(N)}=W^{(N)\dagger}$ at every even length $N>1$. Its source
    ranks equal the physical dimension $d$. In the canonical identifications
    of both rank spaces with $\mathbb C^d$, write these same gates as $u_0,v_0$.
    There exist a unitary $x$ on $\mathbb C^d\otimes\mathbb C^d$ and
    $\sigma\in\{1,-1\}$ such that
    \begin{align}
        (v_0)_{23}(u_0\otimes u_0)
         & =\sigma(x\otimes x^\dagger)(v_0)_{23}^\dagger.
    \end{align}
    This gives the four-original-spin equation of
    \cite[lines~1237--1252]{Cirac2017MPU} for the standard form constructed
    from the chosen reduced canonical source. Its length hypothesis is exactly
    at even standard-form lengths, hence original lengths divisible by four.
    The intermediate virtual comparison retains the possible sign phase
    discussed in \cite{gap:mpu_standard_form_parity_gap}. No gauge or local-equation
    witness is assumed. The actual square standard-form data are supplied
    together with the certificate by relabeling the existing half factors.
    No continuity of any such choice is needed. Arbitrary supplied standard
    forms in an explicitly canonical ambient tensor are handled separately in
    Lemma~\ref{lem:mpu_supplied_forward_local_time_reversal}. The unrestricted
    source proposition remains unformalized. The reduced presentation
    restriction is recorded in \cite{gap:mpu_canonical_form_full_support}.
\end{lemma}

\begin{proof}\leanok
    \uses{lem:mpu_supplied_standard_form_source_cuts,
        lem:mpu_normalized_physical_adjoint_source_factors,
        lem:mpu_even_family_supplied_source_gauges}
    The actual coarse source factors have gates
    $A=v_{23}(u\otimes u)$ and $B=v_{23}$. Adjoint symmetry identifies
    the even periodic families of the standard-form tensor and its adjoint.
    The unblocked virtual comparison supplies a possible sign, which cancels
    in the two-letter source gate. Apply the supplied-factor comparison to
    the actual normalized records. In physical-block coordinates its unitary
    gauges $X,Y$ satisfy $B^\dagger=(X\otimes Y)A$ and
    $A^\dagger=B(Y^\dagger\otimes X^\dagger)$. Taking the adjoint of the
    second equation and substituting the first gives
    $A=(YX\otimes XY)A$. Cancel $A$ to obtain $YX\otimes XY=\Id$.
    Thus $YX=c\Id$ and $XY=c^{-1}\Id$. Unitarity gives $Y=cX^\dagger$,
    whence also $XY=c\Id$ and $c=c^{-1}$. Therefore $c=\pm1$.
    The explicit physical-pair and rank identifications transport this
    relation back to the displayed four-spin equation, with $x=X^\dagger$.
    Only one two-site blocking has been used; no odd-length symmetry follows.
\end{proof}

\begin{lemma}[Standard forms from supplied normalized factors]
    \label{lem:mpu_standard_form_from_normalized_source_factors}
    \lean{MPOTensor.IsMPUCanonicalFormII.exists_twoSiteStandardFormData_sourceFactors}
    \leanok
    Let $U$ be a simple reduced canonical-form-II MPU with recorded weight
    $\rho$. Every normalized supplied source factorization at this weight
    provides actual standard-form data for its own gates on the two-site
    block of $U$. No new compact decomposition is assumed. These are the
    standard forms of \cite[lines~603--648]{Cirac2017MPU}.
\end{lemma}

\begin{proof}\leanok
    \uses{SF, FundamentalMPU}
    Compare the supplied factors with the selected factors of the same tensor.
    The normalized factor comparison gives two unitary rank transformations.
    Apply them to the existing open standard-form contraction. Its gates are
    exactly those of the supplied factors.
\end{proof}

\begin{lemma}[Even-ring converse for the constructed canonical standard form]
    \label{lem:mpu_canonical_even_local_time_reversal_iff}
    \lean{MPOTensor.IsMPUCanonicalFormII.mpo_even_eq_adjoint_of_source_local_certificate, MPOTensor.IsMPUCanonicalFormII.mpo_standardForm_even_eq_adjoint_of_local_certificate, MPOTensor.IsMPUCanonicalFormII.standardForm_even_timeReversal_iff_exists_local_certificate}
    \leanok
    Let $U$ be a simple reduced canonical-form-II tensor, and construct its
    actual standard form $W$ with gates $u,v$. Then $W^{(N)}$ is self-adjoint
    at every even length $N>1$ if and only if $r(U)=\ell(U)=d$ and, in
    these canonical rank identifications, there are a unitary $x$ on
    $\mathbb C^d\otimes\mathbb C^d$ and a sign $\sigma$ satisfying
    \begin{align}
        v_{23}(u\otimes u)
         & =\sigma(x\otimes x^\dagger)v_{23}^\dagger.
    \end{align}
    More generally, a normalized supplied source record of a simple reduced
    canonical tensor, with both source ranks equal to its physical dimension,
    gives self-adjoint periodic operators on even rings whenever its two
    gates satisfy the corresponding sign equation. No extra gate, factor,
    or transform witness is assumed in this implication. This is the
    even-ring contraction in \cite[lines~1214--1236]{Cirac2017MPU}.
    Arbitrary supplied standard-form presentations and transposition remain
    separate; the reduced presentation restriction is recorded in
    \cite{gap:mpu_canonical_form_full_support}.
\end{lemma}

\begin{proof}\leanok
    \uses{lem:mpu_canonical_forward_local_time_reversal,
        lem:mpu_standard_form_from_normalized_source_factors,
        lem:mpu_normalized_physical_adjoint_source_factors, FundamentalMPU}
    The forward implication is the preceding construction. For the converse,
    write $A=v_{23}(u\otimes u)$ and $B=v_{23}$. The local equation
    $A=\sigma(x\otimes x^\dagger)B^\dagger$ gives
    $B^\dagger=(X\otimes Y)A$ and
    $A^\dagger=B(Y^\dagger\otimes X^\dagger)$ with
    $X=x^\dagger$ and $Y=\sigma x$. Both are unitary since $\sigma^2=1$.
    The actual normalized source and adjoint factors provide standard forms
    with gates $A,B$ and $B^\dagger,A^\dagger$, respectively. Close the two
    gate relations on each positive blocked ring by cancelling the internal
    unitaries. These are precisely the even rings of $W$. The explicit
    coordinate dictionary transports the given four-original-spin equation
    to these gate relations without enlarging its local antecedent.
\end{proof}

\begin{lemma}[Admissible local characterization of time reversal]
    \label{prop:mpu_admissible_local_time_reversal}
    \notready
    \discussion{5714}
    % **Scope restriction (two-site blocked comparison):** the tensor obtained by
    % blocking two sites, and the adjoint comparison tensor to which the
    % fundamental theorem is applied, carry reduced full-support source data.
    % No preservation under blocking is asserted. Documented in
    % docs/paper-gaps/mpu_canonical_form_full_support.tex.
    \uses{FundamentalMPU, prop:mpu_admissible_index_well_defined, def:mpu_admissible_index, SF, thm:mpu_canonical_form_ii_stabilization_converse, def:mpu_reduced_full_support_source_datum}
    Let an MPU be in the standard form determined by its given source datum.
    Suppose its one-site tensor
    $\mathcal U$, its two-site block $\mathcal U_2$, and the
    physical-adjoint comparison tensor $\mathcal U_2^\sharp$ are simple and each
    carries its own reduced full-support source datum. Write the one-site
    physical space as $\mathbb C^{d_0}\otimes\mathbb C^{d_0}$. Time-reversal
    symmetry then forces $\operatorname{ind}(\mathcal U)=0$. Under these
    hypotheses, $U^{(N)}=U^{(N)\dagger}$ for every even $N$ if and only if there is a
    unitary $x$ such that
    \begin{align}
        (\Id_1\otimes v_{23}\otimes\Id_4)(u_{12}\otimes u_{34})
         & =\pm(x_{12}\otimes x_{34}^\dagger)
        (\Id_1\otimes v_{23}^\dagger\otimes\Id_4).
    \end{align}
    % Source lines: paper_v2.tex 1201--1252.
\end{lemma}

\begin{proof}
    \uses{FundamentalMPU, thm:mpu_canonical_form_ii_stabilization_converse, prop:mpu_admissible_index_well_defined, def:mpu_admissible_index}
    Under time reversal, $\mathcal U_2$ and $\mathcal U_2^\sharp$ generate the
    same MPU. Applying the same-tensor converse to each given pair produces
    presentations that record the original fixed matrices, and Theorem~\ref{FundamentalMPU}
    then identifies their source spaces,
    so their corresponding source-cut ranks agree. Physical adjunction exchanges
    the two cuts. Hence
    \begin{align}
        r(\mathcal U_2)
         & =r(\mathcal U_2^\sharp)
        =\ell(\mathcal U_2).
    \end{align}
    Definition~\ref{def:mpu_admissible_index} therefore gives
    $\operatorname{ind}(\mathcal U_2)=0$, and blocking independence gives
    $\operatorname{ind}(\mathcal U)=0$. Applying Theorem~\ref{FundamentalMPU}
    to the standard forms of $\mathcal U_2$ and
    $\mathcal U_2^\sharp$ as in the source proof gives the forward implication. The converse follows by
    contracting the displayed local relation around every even periodic chain.
    % Source lines: paper_v2.tex 1201--1252.
\end{proof}

\begin{lemma}[Physical transposition of normalized source factors]
    \label{lem:mpu_physical_transpose_source_factors}
    \lean{MPOTensor.physicalTransposeTensor, MPOTensor.physicalTransposeTensor_apply, MPOTensor.physicalTransposeTensor_physicalTransposeTensor, MPOTensor.mpo_physicalTransposeTensor, MPOTensor.sourceCutM₁_physicalTransposeTensor, MPOTensor.sourceCutM₂_physicalTransposeTensor, MPOTensor.source_ranks_physicalTransposeTensor, MPOTensor.IsMPUCanonicalFormII.physicalTransposeTensor, MPOTensor.IsMPUCanonicalFormII.isMPUSimple_physicalTransposeTensor, MPOTensor.IsMPUCanonicalFormII.sourceRanks_eq_phys_of_even_mpo_eq_transpose, MPOTensor.SourceFactors.physicalTranspose, MPOTensor.SourceFactors.sourceU_physicalTranspose, MPOTensor.SourceFactors.sourceV_physicalTranspose, MPOTensor.IsMPUCanonicalFormII.sourceFactors_physicalTranspose_of_even_mpo_eq_transpose}
    \leanok
    Physical transposition replaces $U^{ij}$ by $U^{ji}$ without reflecting
    its virtual indices. It is an involution, and its periodic operators are
    $(U^{(N)})^T$ at every length. Its source cuts are
    $M_1(\mathcal U^T)=M_2(\mathcal U)^T$ and
    $M_2(\mathcal U^T)=M_1(\mathcal U)^T$, so it exchanges the two source
    ranks. It preserves simplicity and a reduced canonical-form-II
    presentation with its real diagonal weight $\rho$.

    For a simple tensor in this presentation, take any normalized source
    factorization at $\rho$. If $r=\ell=d$, normalized factors for the
    physical transpose are
    \begin{align}
        (X_1',Y_1',Z_1') & =(Y_2^T,X_2^T,\overline{X_2}), \\
        (X_2',Y_2',Z_2') & =(Y_1^T,X_1^T,
        (\Id_d\otimes\rho)\overline{X_1}),
    \end{align}
    in the canonical exchanged-rank coordinates. Their actual gates are
    $u'=v^T$ and $v'=u^T$. Symmetry
    $U^{(N)}=(U^{(N)})^T$ at every even length $N>1$ forces $r=\ell=d$,
    so under this symmetry no square-rank or normalization hypothesis is added
    to the factor construction. These are the physical and factor
    comparisons in \cite[lines~1258--1288]{Cirac2017MPU}.
    This does not establish the local equivalence criterion below. The reduced
    presentation restriction is recorded in
    \cite{gap:mpu_canonical_form_full_support}.
\end{lemma}

\begin{proof}\leanok
    \uses{lem:mpu_normalized_physical_adjoint_source_factors,
        lem:mpu_even_family_supplied_source_gauges}
    Compose physical adjunction with entrywise conjugation. The real diagonal
    weight is fixed under conjugation, so the actual adjoint factor
    normalizations give the displayed transpose factors and gates. Under
    even-chain symmetry the even-family virtual comparison preserves the
    source ranks. Since transposition exchanges them, they agree. Simplicity
    gives $r\ell=d^2$, hence both equal $d$.
\end{proof}

\begin{lemma}[Forward transposition phase for actual canonical source factors]
    \label{lem:mpu_canonical_source_transposition_phase}
    \lean{MPOTensor.IsMPUCanonicalFormII.exists_source_transposition_phase_certificate, MPOTensor.IsMPUCanonicalFormII.exists_source_transposition_exponential_certificate, MPOTensor.IsMPUCanonicalFormII.exists_local_transposition_certificate, MPOTensor.IsMPUCanonicalFormII.exists_local_transposition_exponential_certificate}
    \leanok
    Let $U$ be a simple reduced canonical tensor and construct its actual
    standard form $W$ with gates $u,v$. If
    $W^{(N)}=(W^{(N)})^T$ at every even length $N>1$, then $r(U)=\ell(U)=d$.
    In the corresponding canonical rank identifications there are a unitary
    $x$ on $\mathbb C^d\otimes\mathbb C^d$ and a phase $c$, $|c|=1$, such that
    \begin{align}
        v_{23}(u\otimes u) & =c(x\otimes x^T)v_{23}^T.
    \end{align}
    The scalar can equivalently be written as $c=e^{i\phi}$ with real $\phi$.
    More generally, even-chain transpose symmetry of a simple reduced canonical
    tensor gives the corresponding phase equation for any normalized source
    factors at its specified canonical weight. All comparison gauges are
    derived from the actual periodic family. These are forward consequences
    of \cite[lines~1258--1288]{Cirac2017MPU}, under the presentation restriction
    in \cite{gap:mpu_canonical_form_full_support}. The unrestricted supplied
    standard forms, converse and phase deformation remain separate.
\end{lemma}

\begin{proof}\leanok
    \uses{lem:mpu_physical_transpose_source_factors,
        lem:mpu_even_family_supplied_source_gauges,
        lem:mpu_canonical_forward_local_time_reversal}
    The actual transpose factors have gates $B^T,A^T$. Their two gauge
    equations are $B^T=(X\otimes Y)A$ and
    $A^T=B(Y^\dagger\otimes X^\dagger)$. Transpose the second and substitute
    the first, then cancel $A$, to obtain
    $(\overline YX)\otimes(\overline XY)=\Id$.
    Therefore $\overline YX=c\Id$. Unitarity forces $|c|=1$ and
    $\overline Y=cX^\dagger$. Substitution gives
    $A=c(X^\dagger\otimes\overline X)B^T$. Take $x=X^\dagger$ and use
    the explicit four-spin coordinate dictionary. The principal argument of
    $c$ supplies the pointwise exponential form; no continuity of that choice
    is asserted.
\end{proof}

\begin{lemma}[Even-ring converse for the canonical transposition phase]
    \label{lem:mpu_canonical_even_transposition_iff}
    \lean{MPOTensor.IsMPUCanonicalFormII.mpo_even_eq_transpose_of_source_local_certificate, MPOTensor.IsMPUCanonicalFormII.mpo_standardForm_even_eq_transpose_of_local_certificate, MPOTensor.IsMPUCanonicalFormII.standardForm_even_transposition_iff_exists_local_certificate}
    \leanok
    Let $U$ be a simple reduced canonical tensor, and construct its actual
    standard form $W$ with gates $u,v$. Then $W^{(N)}=(W^{(N)})^T$ at every
    even length $N>1$ if and only if $r(U)=\ell(U)=d$ and, in the canonical
    rank identifications, there are a unitary $x$ on
    $\mathbb C^d\otimes\mathbb C^d$ and a phase $c$, $|c|=1$, satisfying
    \begin{align}
        v_{23}(u\otimes u) & =c(x\otimes x^T)v_{23}^T.
    \end{align}
    A phase is equivalently $e^{i\phi}$ with real $\phi$. More generally,
    the corresponding local phase equation for any normalized source
    factorization of a simple reduced canonical tensor, with square source
    ranks, implies transpose symmetry on its even periodic rings.
    These are the two implications in \cite[lines~1266--1288]{Cirac2017MPU}
    for the constructed canonical source form. Arbitrary supplied standard
    forms remain separate; see
    \cite{gap:mpu_canonical_form_full_support}.
\end{lemma}

\begin{proof}\leanok
    \uses{lem:mpu_canonical_source_transposition_phase,
        lem:mpu_standard_form_from_normalized_source_factors,
        lem:mpu_physical_transpose_source_factors, FundamentalMPU}
    For the converse write $A=c(x\otimes x^T)B^T$ and put
    $X=x^\dagger$, $Y=\overline c\,\overline x$. Since $|c|=1$, both
    matrices are unitary and satisfy $B^T=(X\otimes Y)A$ and
    $A^T=B(Y^\dagger\otimes X^\dagger)$. Use the actual source and transpose
    factor normalizations to construct standard forms with gates $A,B$ and
    $B^T,A^T$, respectively. Closing the two gauge relations cancels each
    internal unitary and gives equality on the even rings of the original
    tensor. The four-spin dictionary transports the given local equation
    without enlarging its antecedent.
\end{proof}

\begin{lemma}[Local transposition for any supplied standard form in a canonical ambient tensor]
    \label{lem:mpu_supplied_even_transposition_criterion}
    \lean{MPOTensor.TwoSiteStandardFormData.intermediateRanks_eq_phys_of_even_mpo_eq_transpose,
        MPOTensor.TwoSiteStandardFormData.exists_transposition_certificate_of_square_ranks,
        MPOTensor.TwoSiteStandardFormData.exists_transposition_certificate,
        MPOTensor.TwoSiteStandardFormData.mpo_even_eq_transpose_of_transposition_certificate,
        MPOTensor.TwoSiteStandardFormData.even_transposition_iff_exists_certificate,
        MPOTensor.TwoSiteStandardFormData.exists_transposition_exponential_certificate,
        Matrix.reindex_transposition_certificate}
    \leanok
    Let $W$ be explicitly in reduced canonical form II and let $u,v$ be
    any supplied standard-form gates on pairs of $d$-dimensional sites.
    Then $W^{(N)}=W^{(N)T}$ for every even $N>1$ if and only if the
    intermediate dimensions satisfy $r=\ell=d$ and, under these coordinate
    identifications, there are a unitary $x$ on
    $\mathbb C^d\otimes\mathbb C^d$ and $|c|=1$ such that
    \[
        v_{23}(u\otimes u)=c(x\otimes x^T)v_{23}^T.
    \]
    The phase can be written $c=e^{i\phi}$ with a pointwise chosen real
    $\phi$. The local equation is preserved under any common coordinate
    relabeling of the pair spaces. No normalization of the supplied half
    factors is required. This is the explicit canonical-ambient,
    even-standard-form-length restriction of
    \cite[lines~1258--1288]{Cirac2017MPU}; the unrestricted source criterion
    retains the presentation restriction in
    \cite{gap:mpu_canonical_form_full_support}.
\end{lemma}
\begin{proof}\leanok
    \uses{lem:mpu_supplied_four_spin_source_comparison,
        lem:mpu_canonical_source_transposition_phase,
        lem:mpu_canonical_even_transposition_iff}
    Even-family comparison with the actual transposed tensor exchanges
    the actual source ranks. The supplied cut-rank identities therefore
    give $r=\ell=d$. Choose the actual normalized source factors and
    compare their gates with $A=v_{23}(u\otimes u)$ and $B=v_{23}$:
    $\widetilde A=(X\otimes Y)A$ and
    $\widetilde B=B(Y^\dagger\otimes X^\dagger)$.
    The normalized transpose certificate with unitary $z$ transports to
    the displayed equation with $x=X^\dagger z\overline Y$.
    Conversely the inverse comparison gives the normalized certificate
    with $z=XxY^T$, whose actual even-ring contraction proves transpose
    symmetry. The argument of the resulting unit complex phase gives
    its displayed exponential form. No further physical blocking is used.
\end{proof}

\begin{lemma}[Continuous removal of the local transposition phase]
    \label{lem:mpu_transposition_phase_deformation}
    \lean{MPOTensor.transpositionPhaseScalar, MPOTensor.norm_transpositionPhaseScalar,
        MPOTensor.transpositionPhaseScalar_sq_mul, MPOTensor.continuous_transpositionPhaseScalar,
        MPOTensor.transpositionPhaseScalar_zero,
        MPOTensor.TwoSiteStandardFormData.transpositionPhaseDeformation,
        MPOTensor.mpo_transpositionPhaseDeformation_eq_transpose,
        MPOTensor.continuous_transpositionPhaseDeformation,
        MPOTensor.transposition_local_phase_deformation,
        MPOTensor.transposition_local_phase_deformation_one}
    \leanok
    Let $W$ have supplied two-site standard-form gates $u,v$ and half factors
    $X_1,X_2$, and let $\phi\in\mathbb R$. Set
    $\delta_t=e^{-it\phi/2}$, $W_t=\delta_t W$ and $u_t=\delta_t u$.
    Then $(u_t,v,X_1,X_2)$ is an actual standard form of $W_t$, and both
    $t\mapsto W_t$ and $t\mapsto u_t$ are continuous, starting at $W,u$.
    At every chain length where $W^{(N)}$ is transpose invariant, so is
    $W_t^{(N)}$. If the square gates satisfy
    $v_{23}(u\otimes u)=e^{i\phi}(x\otimes x^T)v_{23}^T$, then
    \begin{align}
        v_{23}(u_t\otimes u_t)
         & =e^{i(1-t)\phi}(x\otimes x^T)v_{23}^T.
    \end{align}
    In particular the endpoint $t=1$ has local phase one, with the same
    shifted gate and virtual unitary. This is the explicit deformation in
    \cite[lines~1285--1288]{Cirac2017MPU}, without a continuous choice of
    arguments for a varying family of phases.
\end{lemma}

\begin{proof}\leanok
    \uses{SF}
    Since $|\delta_t|=1$, multiplying the unshifted gate by $\delta_t$
    preserves its two unitary identities. The open contraction is linear in
    that gate, giving exactly $W_t$. The periodic operators obey
    $W_t^{(N)}=\delta_t^N W^{(N)}$, which preserves transpose symmetry.
    The two unshifted gates in the local equation contribute $\delta_t^2$,
    and $\delta_t^2e^{i\phi}=e^{i(1-t)\phi}$. The complex exponential and
    scalar multiplication give continuity.
\end{proof}

\begin{lemma}[Local characterization of transposition]
    \label{prop:local-char-transposition}
    \notready
    \discussion{5714}
    \uses{FundamentalMPU, SF, def:index}
    Let an MPU be in standard form, with one-site physical space
    $\mathbb C^{d_0}\otimes\mathbb C^{d_0}$. Transposition symmetry forces
    $\operatorname{ind}(\mathcal U)=0$. The MPU satisfies
    $U^{(N)}=U^{(N)T}$ for every even $N$ if and only if there are a unitary
    $x$ and a phase $e^{i\phi}$ such that
    \begin{align}
        (\Id_1\otimes v_{23}\otimes\Id_4)(u_{12}\otimes u_{34})
         & =e^{i\phi}(x_{12}\otimes x_{34}^T)
        (\Id_1\otimes v_{23}^T\otimes\Id_4).
    \end{align}
    % Source lines: paper_v2.tex 1257--1288.
\end{lemma}

\begin{proof}
    \uses{FundamentalMPU, index-well-defined, def:index,
        def:mpu_source_matrix_one, def:mpu_source_matrix_two,
        def:mpu_right_rank, def:mpu_left_rank}
    Under transposition, the two-site block $\mathcal U_2$ and its transposed
    comparison tensor $\mathcal U_2^{\mathsf T}$ generate the same MPU. The
    fundamental theorem identifies their source spaces, so their corresponding
    source-cut ranks agree. Swapping the physical input and output indices
    transposes the first source matrix into the second and therefore exchanges
    the two ranks. Hence
    \begin{align}
        r(\mathcal U_2)
         & =r(\mathcal U_2^{\mathsf T})
        =\ell(\mathcal U_2).
    \end{align}
    Thus $\operatorname{ind}(\mathcal U_2)=0$ by
    Definition~\ref{def:index}, and blocking independence gives
    $\operatorname{ind}(\mathcal U)=0$.

    Applying the two gauge equations of the fundamental theorem to the standard
    forms of $\mathcal U_2$ and $\mathcal U_2^{\mathsf T}$ shows that their compatibility
    gives $y=e^{i\phi}x^T$ for some phase $e^{i\phi}$, which yields the
    displayed local relation. Conversely, contract that relation around an
    even periodic chain; the phases cancel pairwise.
    % Source lines: paper_v2.tex 1257--1288.
\end{proof}

\begin{lemma}[Admissible local characterization of transposition]
    \label{prop:mpu_admissible_local_transposition}
    \notready
    \discussion{5714}
    % **Scope restriction (two-site blocked comparison):** the tensor obtained by
    % blocking two sites, and the transposed comparison tensor to which the
    % fundamental theorem is applied, carry reduced full-support source data.
    % No preservation under blocking is asserted. Documented in
    % docs/paper-gaps/mpu_canonical_form_full_support.tex.
    \uses{FundamentalMPU, prop:mpu_admissible_index_well_defined, def:mpu_admissible_index, SF, thm:mpu_canonical_form_ii_stabilization_converse, def:mpu_reduced_full_support_source_datum}
    Let an MPU be in the standard form determined by its given source datum.
    Suppose its one-site tensor
    $\mathcal U$, its two-site block $\mathcal U_2$, and the transposed
    comparison tensor $\mathcal U_2^{\mathsf T}$ are simple and each carries its
    own reduced full-support source datum. Write the one-site physical space as
    $\mathbb C^{d_0}\otimes\mathbb C^{d_0}$. Transposition symmetry then forces
    $\operatorname{ind}(\mathcal U)=0$. Under these hypotheses,
    $U^{(N)}=U^{(N)T}$ for every even $N$ if and only if there are a unitary
    $x$ and a phase $e^{i\phi}$ such that
    \begin{align}
        (\Id_1\otimes v_{23}\otimes\Id_4)(u_{12}\otimes u_{34})
         & =e^{i\phi}(x_{12}\otimes x_{34}^T)
        (\Id_1\otimes v_{23}^T\otimes\Id_4).
    \end{align}
    % Source lines: paper_v2.tex 1257--1288.
\end{lemma}

\begin{proof}
    \uses{FundamentalMPU, thm:mpu_canonical_form_ii_stabilization_converse, prop:mpu_admissible_index_well_defined, def:mpu_admissible_index}
    Under transposition, $\mathcal U_2$ and $\mathcal U_2^{\mathsf T}$ generate
    the same MPU. Applying the same-tensor converse to each given pair produces
    presentations that record the original fixed matrices, and Theorem~\ref{FundamentalMPU}
    then identifies their source spaces, so their corresponding source-cut
    ranks agree. Transposition
    exchanges the two cuts. Hence
    \begin{align}
        r(\mathcal U_2)
         & =r(\mathcal U_2^{\mathsf T})
        =\ell(\mathcal U_2).
    \end{align}
    Thus $\operatorname{ind}(\mathcal U_2)=0$ by
    Definition~\ref{def:mpu_admissible_index}, and blocking independence gives
    $\operatorname{ind}(\mathcal U)=0$. Applying Theorem~\ref{FundamentalMPU}
    to the standard forms of $\mathcal U_2$ and
    $\mathcal U_2^{\mathsf T}$ and repeating the time-reversal argument while retaining
    the unrestricted phase $e^{i\phi}$. The converse follows by contracting the
    local relation around every even periodic chain.
    % Source lines: paper_v2.tex 1257--1288.
\end{proof}

\subsection{Equivalence under symmetries}

\begin{definition}[Finite-chain operator symmetry]
    \label{def:mpu_finite_chain_operator_symmetry}
    \lean{MPOTensor.FiniteChainOperatorSymmetry,
        MPOTensor.IsInvariantUnderSymmetry}
    \leanok
    At physical dimension $d$, a finite-chain operator symmetry $\mathcal S$
    specifies a set $A$ of applicable chain lengths and, for every length $N$,
    a map
    \begin{align}
        \mathcal S_{d,N}\colon
        \operatorname{Mat}_{d^N}(\mathbb C)
         & \longrightarrow \operatorname{Mat}_{d^N}(\mathbb C).
    \end{align}
    An MPO tensor $\mathcal W$ is invariant under $\mathcal S$ if
    \begin{align}
        \mathcal S_{d,N}[W^{(N)}] & =W^{(N)}
    \end{align}
    for every $N\in A$.
    % Source lines: paper_v2.tex 1345--1354.
\end{definition}

\begin{definition}[Coherent symmetry family under identity ancillas]
    \label{def:mpu_coherent_symmetry_family}
    \lean{MPOTensor.tensorPhysicalIdOperator,
        MPOTensor.FiniteChainOperatorSymmetryFamily,
        MPOTensor.FiniteChainOperatorSymmetryFamily.at}
    \leanok
    \uses{def:mpu_finite_chain_operator_symmetry, def:equivalent-tensors}
    A coherent symmetry family specifies an action $\mathcal S_{d,N}$ for every
    physical dimension $d$, with one common set $A$ of applicable lengths. For
    an ancilla dimension $x$, let
    \begin{align}
        \iota_x(X)
         & = R_{d,x,N}\bigl(X\otimes I_{x^N}\bigr)R_{d,x,N}^{-1},
    \end{align}
    where $R_{d,x,N}$ is the canonical sitewise identification of the
    $d$-dimensional and $x$-dimensional configurations. Coherence means
    \begin{align}
        \iota_x\bigl(\mathcal S_{d,N}(X)\bigr)
         & = \mathcal S_{dx,N}\bigl(\iota_x(X)\bigr).
    \end{align}
    Thus the action on the full enlarged operator algebra is part of the given data; it
    is not obtained by extending its values on operators $X\otimes I$.
    % Source lines: paper_v2.tex 1356--1364.
\end{definition}

\begin{lemma}[Identity-ancilla invariance under a symmetry family]
    \label{thm:mpu_symmetry_identity_ancilla_invariant}
    \lean{MPOTensor.IsInvariantUnderSymmetry.tensorPhysicalId}
    \leanok
    \uses{def:mpu_coherent_symmetry_family, def:equivalent-tensors}
    If $\mathcal U$ is invariant under the $d$-dimensional component of a
    coherent symmetry family, then $\mathcal U^{(x)}$ is invariant under its
    $dx$-dimensional component.
\end{lemma}

\begin{proof}\leanok
    The finite-chain operator of $\mathcal U^{(x)}$ is $\iota_x(U^{(N)})$.
    Semiconjugacy and invariance of $U^{(N)}$ give
    \begin{align}
        \mathcal S_{dx,N}\bigl(\iota_x(U^{(N)})\bigr)
         & = \iota_x\bigl(\mathcal S_{d,N}(U^{(N)})\bigr)
        = \iota_x(U^{(N)}).
    \end{align}
\end{proof}

\begin{definition}[Blocked finite-chain symmetry]
    \label{def:mpu_blocked_finite_chain_symmetry}
    \lean{MPOTensor.FiniteChainOperatorSymmetry.block,
        MPOTensor.FiniteChainOperatorSymmetry.block_action}
    \leanok
    \uses{def:mpu_finite_chain_operator_symmetry, lem:eval_word_block}
    For a blocking length $k$, the blocked symmetry has applicable lengths
    \begin{align}
        A_k & = \{N\mid Nk\in A\}.
    \end{align}
    Its action is obtained by conjugating $\mathcal S_{d,Nk}$ along the
    standard identification
    $(\mathbb C^{d^k})^{\otimes N}\cong(\mathbb C^d)^{\otimes Nk}$.
    % Source line: paper_v2.tex 1366.
\end{definition}

\begin{lemma}[Preservation of symmetry invariance under blocking]
    \label{thm:mpu_symmetry_blocking_invariant}
    \lean{MPOTensor.IsInvariantUnderSymmetry.blockTensor}
    \leanok
    \uses{def:mpu_blocked_finite_chain_symmetry, thm:mpdo_closed_chain_physical_blocking}
    If $\mathcal U$ is invariant under $\mathcal S$, then its $k$-site blocking
    is invariant under the blocked symmetry.
\end{lemma}

\begin{proof}\leanok
    Let $R_{d,N,k}$ be the standard identification from blocked length-$N$
    configurations to unblocked length-$Nk$ configurations. For $N\in A_k$,
    one has $Nk\in A$ and
    \begin{align}
        \mathcal S_{k,N}\bigl[U_k^{(N)}\bigr]
         & = R_{d,N,k}^{-1}\!\left(
        \mathcal S_{d,Nk}\!\left[
                R_{d,N,k}\bigl(U_k^{(N)}\bigr)
        \right]\right)                               \\
         & = R_{d,N,k}^{-1}\!\left(
        \mathcal S_{d,Nk}\bigl[U^{(Nk)}\bigr]\right) \\
         & = R_{d,N,k}^{-1}\bigl(U^{(Nk)}\bigr)      \\
         & = U_k^{(N)}.
    \end{align}
    The second and final equalities are the inverse reindexing identities, and
    the third is invariance of $U^{(Nk)}$ because $Nk\in A$.
\end{proof}

\begin{definition}[Strict equivalence under a symmetry]
    \label{def:strictly-equivalent-symmetry}
    \lean{MPOTensor.StrictlyEquivalentUnderSymmetry,
        MPOTensor.StrictlyEquivalentUnderSymmetry.of_fixedBond}
    \leanok
    \discussion{6948}
    \uses{def:strictly-equivalent-tensors, def:mpu_finite_chain_operator_symmetry}
    % **Local fix:** paper_v2.tex 1345--1353 repeats the symmetry-transformed
    % operator on the right-hand side. Invariance requires the original MPU.
    % Documented in docs/paper-gaps/mpu_canonical_form_full_support.tex.
    Two canonical-form tensors $\mathcal U$
    and $\mathcal V$ are strictly equivalent under a symmetry $\mathcal S$ if
    they have the same physical dimension and, for some common bond dimension
    $D'$ as in Definition~\ref{def:strictly-equivalent-tensors}, there is a
    continuous path $\mathcal W(p)$ of MPU tensors of bond dimension $D'$ from
    $\mathcal U_{\uparrow D'}$ to $\mathcal V_{\uparrow D'}$,
    not necessarily in canonical form, such that
    \begin{align}
        \mathcal S[W^{(N)}(p)] & =W^{(N)}(p)
    \end{align}
    for every $p\in[0,1]$ and every applicable chain length $N$. When
    $\mathcal U$ and $\mathcal V$ have one bond dimension $D$, such a path of
    bond dimension $D$ from $\mathcal U$ to $\mathcal V$ realizes the relation
    with $D'=D$.
    % Source lines: paper_v2.tex 1340--1354.
\end{definition}

\begin{lemma}[Forgetting symmetry from strict equivalence]
    \label{lem:mpu_strict_symmetry_equivalence_forget}
    \lean{MPOTensor.StrictlyEquivalentUnderSymmetry.toStrictlyEquivalent}
    \leanok
    \uses{def:strictly-equivalent-symmetry, def:strictly-equivalent-tensors}
    Strict equivalence under a symmetry implies strict equivalence.
\end{lemma}

\begin{proof}\leanok
    The same path lies in the larger locus obtained by forgetting invariance
    under $\mathcal S$.
\end{proof}

\begin{definition}[Admissible strict equivalence under a symmetry]
    \label{def:mpu_admissible_strict_symmetry_equivalence}
    \notready
    \discussion{6017}
    \uses{def:strictly-equivalent-tensors, def:mpu_admissible_equivalence_datum}
    Two tensors $\mathcal U$ and $\mathcal V$ are strictly equivalent under a
    symmetry $\mathcal S$ if they have the same physical dimension and there is
    a continuous path $\mathcal W(p)$ of MPU tensors from $\mathcal U$ to
    $\mathcal V$ such that
    \begin{align}
        \mathcal S[W^{(N)}(p)] & =W^{(N)}(p)
    \end{align}
    for every $p\in[0,1]$ and every chain length $N$ under consideration. The
    path need not be in canonical form. The equivalence is admissible if its
    connecting path carries an admissible equivalence datum for the
    blocked path tensor and the fixed conjugate, physical-adjoint, transposed,
    and $\mathcal S$-image comparison paths, and if the common blocking length
    $m$ of that datum satisfies $m\leq D^4$.
    % Source lines: paper_v2.tex 1340--1354.
\end{definition}

\begin{definition}[Equivalence under a symmetry]
    \label{def:equivalent-symmetry}
    \notready
    \discussion{7422}
    \uses{def:strictly-equivalent-symmetry, def:equivalent-tensors}
    Let $\mathcal U$ and $\mathcal V$ have physical dimensions $d_a$ and
    $d_b$. They are equivalent under a symmetry $\mathcal S$ if there are a
    positive blocking length $k$ and coprime positive ancilla dimensions
    $p_a,p_b$ with $p_ad_a=p_bd_b$ such that $(\mathcal U^{(p_a)})_k$ and
    $(\mathcal V^{(p_b)})_k$ are strictly equivalent under $\mathcal S$.
    % Source lines: paper_v2.tex 1356--1366.
\end{definition}

\begin{definition}[Coherent-family equivalence under a symmetry]
    \label{def:mpu_coherent_family_equivalent_symmetry}
    \lean{MPOTensor.EquivalentUnderSymmetry}
    \leanok
    \discussion{7015}
    \uses{def:strictly-equivalent-symmetry, def:equivalent-tensors,
        def:mpu_coherent_symmetry_family, def:mpu_blocked_finite_chain_symmetry}
    % **Scope restriction (coherent family):** the symmetry is
    % supplied at every physical dimension and semiconjugates the canonical
    % identity-ancilla embeddings. Documented in
    % docs/paper-gaps/mpu_symmetry_ancilla_transport.tex.
    Let $\mathcal U$ and $\mathcal V$
    have physical dimensions $d_a$ and $d_b$. They are coherently equivalent
    under a dimension-indexed symmetry family $\mathcal S$ if they are MPUs and
    there are a positive blocking length $k$ and coprime positive ancilla
    dimensions $p_a,p_b$ with $p_ad_a=p_bd_b$ such that
    $(\mathcal U^{(p_a)})_k$ and
    $(\mathcal V^{(p_b)})_k$ are strictly equivalent under the $k$-site
    blocking of the given enlarged-dimension symmetry. The family obeys the
    identity-ancilla semiconjugacy law, and its applicable blocked lengths are
    precisely the lengths $N$ for which $Nk\in A$. The ancillas are attached
    before the common blocking.
    % Source lines: paper_v2.tex 1356--1366, under the stated restrictions.
\end{definition}

\begin{lemma}[Forgetting symmetry from coherent stabilized equivalence]
    \label{thm:mpu_symmetry_equivalence_forget}
    \lean{MPOTensor.EquivalentUnderSymmetry.toEquivalent}
    \leanok
    \uses{def:mpu_coherent_family_equivalent_symmetry, def:equivalent-tensors,
        lem:mpu_strict_symmetry_equivalence_forget}
    Equivalence under a coherent symmetry family implies ordinary
    equivalence after the same ancillas and blocking.
\end{lemma}

\begin{proof}\leanok
    Retain the MPU hypotheses, positive coprime ancilla dimensions, common
    positive blocking length, and physical-dimension equality from the
    symmetry-preserving equivalence data. Forgetting invariance from its final strict
    path gives the required ordinary stabilized equivalence.
\end{proof}

\begin{definition}[Admissible equivalence under a symmetry]
    \label{def:mpu_admissible_symmetry_equivalence}
    \notready
    \discussion{6017}
    \uses{def:mpu_admissible_strict_symmetry_equivalence, def:equivalent-tensors, def:mpu_admissible_equivalence_datum}
    Let $\mathcal U$ and $\mathcal V$ have physical dimensions $d_a$ and
    $d_b$. They are equivalent under a symmetry $\mathcal S$ if there are a
    positive blocking length $k$ and coprime positive ancilla dimensions
    $p_a,p_b$ with $p_ad_a=p_bd_b$ such that $(\mathcal U^{(p_a)})_k$ and
    $(\mathcal V^{(p_b)})_k$ are strictly equivalent under $\mathcal S$. Thus
    the ancillas are attached before the common blocking. The equivalence is
    admissible if this strict-equivalence path carries an admissible
    equivalence datum for the blocked path itself and the fixed conjugate,
    physical-adjoint, transposed, and $\mathcal S$-image comparison paths, with
    common blocking length at most $D^4$.
    % Source lines: paper_v2.tex 1356--1364.
\end{definition}

\begin{lemma}[Standard-form path criterion]
    \label{lem:equivalent-symmetry}
    \notready
    \discussion{6017}
    \uses{def:strictly-equivalent-symmetry, coruvSF, FundamentalMPU}
    Two MPUs described by simple tensors are strictly equivalent under
    $\mathcal S$ if and only if their standard-form unitaries are joined by
    continuous families $\widetilde u(p)$ and $\widetilde v(p)$ whose generated
    MPU is invariant under $\mathcal S$ for every $p\in[0,1]$.
    % Source lines: paper_v2.tex 1366--1380.
\end{lemma}

\begin{lemma}[Admissible standard-form path criterion]
    \label{lem:mpu_admissible_symmetry_path_criterion}
    \notready
    \discussion{6017}
    % **Scope restriction (admissible symmetry path):** the strict-equivalence
    % witness carries the path datum of
    % Definition~\ref{def:mpu_admissible_equivalence_datum}. Documented in
    % docs/paper-gaps/mpu_canonical_form_full_support.tex.
    \uses{def:mpu_admissible_strict_symmetry_equivalence, cor:mpu_admissible_continuous_standard_form, FundamentalMPU, thm:mpu_canonical_form_ii_stabilization_converse, def:mpu_admissible_equivalence_datum}
    Fix a positive integer $m\leq D^4$. The $m$-site blocked endpoint
    MPUs $U_1^{(mN)}$ and $U_2^{(mN)}$ are admissibly strictly equivalent under
    $\mathcal S$ if and only if there are continuous families of unitaries
    $\widetilde u(p)$ and $\widetilde v(p)$ whose endpoints generate
    $U_1^{(mN)}$ and $U_2^{(mN)}$, whose generated MPU is invariant under
    $\mathcal S$ for every $p\in[0,1]$, and whose tensor path and fixed
    conjugate, physical-adjoint, transposed, and $\mathcal S$-image comparison
    paths carry continuously varying reduced full-support source data. These
    families arise from the blocked path $\mathcal W(p)_m$ in an
    admissible equivalence datum of common length $m$; the blocked path itself is
    the length-one admissible path for the displayed endpoint families,
    matching in \cite[Lemma following the symmetry-equivalence definitions]{Cirac2017MPU} the blocking
    lemma.
    % Source lines: paper_v2.tex 1366--1380.
\end{lemma}

\begin{lemma}[Conjugation symmetry when one source rank is odd]
    \label{conjphs}
    \notready
    \discussion{5713}
    \uses{prop1conj, corconj, lem:equivalent-symmetry}
    Suppose that either $r$ or $\ell$ is odd. Two conjugation-invariant MPUs
    $\mathcal U_1$ and $\mathcal U_2$ are strictly equivalent under
    conjugation if and only if
    \begin{align}
        \det(u_1v_1) & =\det(u_2v_2).
    \end{align}
    % Source lines: paper_v2.tex 1401--1446.
\end{lemma}

\begin{lemma}[Admissible conjugation symmetry when one source rank is odd]
    \label{thm:mpu_admissible_conjugation_odd_rank}
    \notready
    \discussion{5713}
    % **Scope restriction (endpoints and admissible symmetry path):** each
    % original endpoint and its conjugate comparison carry source data, and
    % equivalence means admissible strict equivalence in
    % Definition~\ref{def:mpu_admissible_equivalence_datum}. Documented in
    % docs/paper-gaps/mpu_canonical_form_full_support.tex.
    \uses{prop:mpu_admissible_conjugation_standard_forms, lem:mpu_admissible_symmetry_path_criterion, def:mpu_admissible_equivalence_datum, def:mpu_reduced_full_support_source_datum}
    For $i=1,2$, suppose the original one-site endpoint tensor
    $\mathcal U_i$ and its conjugate comparison tensor are simple and each
    carries reduced full-support source data. Suppose that either $r$ or $\ell$
    is odd and that the standard-form interpolation determined below carries an
    admissible path datum of blocking length one. Then $\mathcal U_1$ and
    $\mathcal U_2$ are admissibly strictly equivalent
    under conjugation symmetry if and only if
    \begin{align}
        \det(u_1v_1) & =\det(u_2v_2).
    \end{align}
    % Source lines: paper_v2.tex 1401--1446.
\end{lemma}

\begin{corollary}[Equivalence of conjugation-symmetric MPUs after blocking]
    \label{cor:all-equivalent-under-conjugation}
    \notready
    \uses{conjphs, def:equivalent-symmetry, lem:mpu_blocking}
    If either $r$ or $\ell$ is odd, all conjugation-invariant MPUs with these
    source ranks are equivalent under conjugation symmetry.
    % Source lines: paper_v2.tex 1450--1460.
\end{corollary}

\begin{corollary}[Admissible equivalence of conjugation-symmetric MPUs after blocking]
    \label{cor:mpu_admissible_all_conjugation_equivalent}
    \notready
    \discussion{5713}
    % **Scope restriction (endpoints, blocking, and admissible path):** the
    % one-site and two-site endpoints and their conjugate comparisons carry
    % source data, and the resulting symmetry path is admissible. No
    % preservation under blocking is asserted. Documented in
    % docs/paper-gaps/mpu_canonical_form_full_support.tex.
    \uses{prop:mpu_admissible_conjugation_standard_forms, cor:mpu_admissible_equiv_under_conjugation, def:mpu_admissible_symmetry_equivalence, def:mpu_admissible_equivalence_datum, def:mpu_reduced_full_support_source_datum}
    For $i=1,2$, suppose the endpoint tensor $\mathcal U_i$, its conjugate
    comparison $\overline{\mathcal U_i}$, the two-site block
    $\mathcal U_{i,2}$, and its conjugate comparison
    $\overline{\mathcal U_{i,2}}$ are simple and each carries reduced
    full-support source data. Suppose that either $r$ or $\ell$ is odd and that
    the resulting standard-form interpolation carries an admissible path datum
    of blocking length one. Then
    $\mathcal U_1$ and $\mathcal U_2$ are admissibly equivalent under
    conjugation symmetry.
    % Source lines: paper_v2.tex 1450--1460.
\end{corollary}

\begin{proof}
    \uses{prop:mpu_admissible_conjugation_standard_forms, cor:mpu_admissible_equiv_under_conjugation}
    If either $r$ or $\ell$ is odd, both endpoints lie in Case I of
    Lemma~\ref{prop:mpu_admissible_conjugation_standard_forms}. The full
    admissible classification in
    Lemma~\ref{cor:mpu_admissible_equiv_under_conjugation} handles both
    possible parities of the source ranks after blocking and gives the claimed
    admissible equivalence.
\end{proof}

\begin{lemma}[Determinants of orthogonal Kronecker products]
    \label{thm:mpu_orthogonal_kronecker_determinant}
    \lean{Matrix.det_eq_one_or_neg_one_of_mem_orthogonalGroup,
        Matrix.det_pow_eq_one_of_mem_orthogonalGroup_of_even,
        Matrix.det_kronecker_eq_one_of_mem_orthogonalGroup}
    \leanok
    Let $O$ and $O'$ be orthogonal matrices over a commutative ring,
    indexed by finite sets of cardinalities $m$ and $n$. For every even
    integer $k\geq0$, $\det(O)^k=1$. Over an integral domain one also has
    $\det(O)\in\{1,-1\}$. If $m$ and $n$ are even, then
    \begin{align}
        \det(O\otimes O') & =\det(O)^n\det(O')^m=1.
    \end{align}
    In \cite{Cirac2017MPU}, these identities give the orthogonal determinant
    calculation.
    % Source: paper_v2.tex lines 1464--1477.
\end{lemma}

\begin{proof}\leanok
    Taking determinants in $OO^T=\Id$ gives $\det(O)^2=1$, and the same holds
    for $O'$. The determinant of a Kronecker product is
    $\det(O)^n\det(O')^m$; both powers are one when $m$ and $n$ are even.
\end{proof}

\begin{lemma}[Determinant of a symplectic Kronecker product]
    \label{thm:mpu_symplectic_kronecker_determinant}
    \lean{Matrix.det_kronecker_eq_one_of_mem_symplecticGroup}
    \leanok
    Let $S$ and $S'$ preserve the standard symplectic forms on spaces of
    dimensions $2p$ and $2q$ over a commutative ring. Then
    \begin{align}
        \det(S\otimes S') & =\det(S)^{2q}\det(S')^{2p}=1.
    \end{align}
    No positivity assumption on either dimension is needed. This theorem
    uses the standard symplectic form; the interleaved form in
    \cite{Cirac2017MPU} requires a separate change of coordinates before
    this calculation applies to its matrices.
    % Source: paper_v2.tex lines 1464--1477.
\end{lemma}

\begin{proof}\leanok
    Since each symplectic determinant equals one, the Kronecker determinant
    formula gives both identities.
\end{proof}

\begin{lemma}[Strict conjugation equivalence for even source ranks]
    \label{prop:conjsym-strict-equiv}
    \notready
    \discussion{5713}
    \uses{prop1conj, corconj, cor:sympl-to-orthogonal, lem:equivalent-symmetry, lemma:conjclass-normalform-continuous}
    Suppose $r$ and $\ell$ are even. Two conjugation-invariant MPUs in the
    standard forms of Lemma~\ref{corconj} are strictly equivalent under
    conjugation if and only if they belong to the same case and
    \begin{align}
        \det(u_1) & =\det(u_2), &
        \det(v_1) & =\det(v_2).
    \end{align}
    % Source lines: paper_v2.tex 1464--1546.
\end{lemma}

\begin{lemma}[Admissible strict conjugation equivalence for even source ranks]
    \label{thm:mpu_admissible_conjugation_even_rank}
    \notready
    \discussion{5713}
    % **Scope restriction (admissible symmetry path):** equivalence means
    % admissible strict equivalence in
    % Definition~\ref{def:mpu_admissible_equivalence_datum}. Documented in
    % docs/paper-gaps/mpu_canonical_form_full_support.tex.
    \uses{prop:mpu_admissible_conjugation_standard_forms, cor:mpu_admissible_symplectic_to_orthogonal, lem:mpu_admissible_symmetry_path_criterion, lemma:conjclass-normalform-continuous, def:mpu_admissible_equivalence_datum}
    Suppose that $r$ and $\ell$ are both even and that the standard-form
    interpolation determined below carries an admissible path datum of blocking
    length one. Choose standard forms from the given source data of two
    conjugation-invariant
    MPUs as in
    Lemma~\ref{prop:mpu_admissible_conjugation_standard_forms}. They are admissibly strictly equivalent under
    conjugation symmetry if and only if they belong
    to the same case of Lemma~\ref{prop:mpu_admissible_conjugation_standard_forms} and
    \begin{align}
        \det(u_1) & =\det(u_2), &
        \det(v_1) & =\det(v_2).
    \end{align}
    % Source lines: paper_v2.tex 1464--1546.
\end{lemma}

\begin{lemma}[Equivalence under conjugation]
    \label{cor:equiv-under-conjugation}
    \notready
    \discussion{5713}
    \uses{corconj, prop:conjsym-strict-equiv, def:equivalent-symmetry}
    Two conjugation-invariant MPUs are equivalent under conjugation symmetry if
    and only if they belong to the same case of Lemma~\ref{corconj}.
    % Source lines: paper_v2.tex 1548--1567.
\end{lemma}

\begin{lemma}[Admissible equivalence under conjugation]
    \label{cor:mpu_admissible_equiv_under_conjugation}
    \notready
    \discussion{5713}
    % **Scope restriction (endpoints, blocking, and admissible path):** the
    % original one-site endpoints and their conjugate comparisons carry source
    % data, all blocked endpoints used below carry their own data, and the
    % symmetry path is admissible. Documented in
    % docs/paper-gaps/mpu_canonical_form_full_support.tex.
    \uses{thm:mpu_admissible_conjugation_odd_rank, thm:mpu_admissible_conjugation_even_rank, prop:mpu_admissible_conjugation_standard_forms, prop:mpu_admissible_index_well_defined, def:mpu_admissible_symmetry_equivalence, def:mpu_admissible_equivalence_datum, def:mpu_reduced_full_support_source_datum}
    For $i=1,2$, suppose the original one-site endpoint tensor
    $\mathcal U_i$, its conjugate comparison $\overline{\mathcal U_i}$, the
    two-site block $\mathcal U_{i,2}$, and its conjugate comparison
    $\overline{\mathcal U_{i,2}}$ are simple and each carries reduced
    full-support source data. Suppose also that the resulting standard-form
    interpolation carries an admissible path datum of blocking length one. Then
    $\mathcal U_1$ and $\mathcal U_2$ are admissibly equivalent under
    conjugation symmetry if and only if they belong to the same case of
    Lemma~\ref{prop:mpu_admissible_conjugation_standard_forms}.
    % Source lines: paper_v2.tex 1548--1567.
\end{lemma}

\begin{proof}
    \uses{thm:mpu_admissible_conjugation_odd_rank, thm:mpu_admissible_conjugation_even_rank, prop:mpu_admissible_conjugation_standard_forms, prop:mpu_admissible_index_well_defined, def:mpu_admissible_symmetry_equivalence}
    Suppose first that either $r$ or $\ell$ is odd. Then both original
    endpoints lie in Case I. After one blocking, write $r^{(2)}$ and
    $\ell^{(2)}$ for the blocked source ranks. Admissible blocking independence,
    applied to the given one-site and two-site endpoint data, gives
    \begin{align}
        r^{(2)} & =dr,
    \end{align}
    and
    \begin{align}
        \ell^{(2)} & =d\ell.
    \end{align}
    For each endpoint,
    \begin{align}
        \det(u_i^{(2)}v_i^{(2)}) & =\det(u_iv_i)^{2r\ell}=1.
    \end{align}
    Blocking preserves Case I. If either $r^{(2)}$ or $\ell^{(2)}$ is odd, the
    odd-rank theorem applied to the blocked endpoints gives admissible strict
    equivalence. If both blocked ranks are even, then $d$ is even because one
    original source rank is odd. Hence $r\ell=d^2$ is even.
    % **Local fix:** paper_v2.tex 1554--1559 states the two separate determinant
    % identities before checking parity. They are used here only when the
    % blocked source ranks are both even, which forces the required even power.
    % Documented in docs/paper-gaps/mpu_canonical_form_full_support.tex.
    Consequently,
    \begin{align}
        \det(u_i^{(2)}) & =(\det(u_i)^2\det(v_i))^{r\ell}=1,
    \end{align}
    and
    \begin{align}
        \det(v_i^{(2)}) & =\det(v_i)^{r\ell}=1.
    \end{align}
    The even-rank theorem now applies to the preserved Case I endpoints. Thus
    the odd-source-rank branch has a single admissible equivalence class without
    imposing a parity condition on the blocked ranks.

    If $r$ and $\ell$ are both even, admissible blocking independence shows that
    the blocked ranks remain even. Since $r\ell$ is even, the same calculations
    give
    \begin{align}
        \det(u_i^{(2)}) & =(\det(u_i)^2\det(v_i))^{r\ell}=1,
    \end{align}
    and
    \begin{align}
        \det(v_i^{(2)}) & =\det(v_i)^{r\ell}=1.
    \end{align}
    Blocking preserves Case I and Case II. The even-rank theorem therefore shows
    that the blocked endpoints are admissibly strictly equivalent exactly when
    they belong to the same case. This is precisely admissible equivalence of
    the original endpoints.
\end{proof}

\begin{lemma}[SWAP trace contraction]
    \label{lem:matrix_trace_swap_kronecker}
    \lean{Matrix.trace_swapMatrix_mul_kronecker,
        Matrix.trace_swapMatrix_mul_kronecker_conjTranspose_of_mem_unitaryGroup}
    \leanok
    For matrices $A,B\in\mathbb C^{n\times n}$ and the exchange matrix $\mathbb S$ on
    $\mathbb C^n\otimes\mathbb C^n$,
    \begin{align}
        \tr\!\left[\mathbb S(A\otimes B)\right] & =\tr(AB).
    \end{align}
    In particular, if $U$ is unitary, then
    \begin{align}
        \tr\!\left[\mathbb S(U\otimes U^\dagger)\right] & =n.
    \end{align}
    % Source identity: paper_v2.tex lines 1608--1610.
\end{lemma}

\begin{proof}\leanok
    In product coordinates, the exchange matrix has entries
    $\mathbb S_{(i,j),(k,l)}=\delta_{i,l}\delta_{j,k}$. Therefore
    \begin{align}
        \tr\!\left[\mathbb S(A\otimes B)\right]
         & =\sum_{i,j}A_{j,i}B_{i,j}
        =\sum_{i,j}A_{i,j}B_{j,i}
        =\tr(AB).
    \end{align}
    Taking $A=U$ and $B=U^\dagger$ gives
    $\tr(UU^\dagger)=\tr(\Id)=n$.
\end{proof}

\begin{lemma}[The exchange intertwines the Kronecker factors]
    \label{thm:matrix_swap_kronecker_intertwine}
    \lean{Matrix.swapMatrix_mul_kronecker, Matrix.kronecker_mul_swapMatrix}
    \leanok
    For matrices $A,B\in\mathbb C^{n\times n}$ and the exchange matrix
    $\mathbb S$ on $\mathbb C^n\otimes\mathbb C^n$,
    \begin{align}
        \mathbb S(A\otimes B) & =(B\otimes A)\mathbb S,
        \label{eq:matrix_swap_kron_left}                \\
        (A\otimes B)\mathbb S & =\mathbb S(B\otimes A).
        \label{eq:matrix_swap_kron_right}
    \end{align}
    The equivalent two-sided exchange identity is stated in
    \cite[lines~1316--1319]{Cirac2017MPU}; the gauge-invariance argument at
    \cite[lines~1624--1630]{Cirac2017MPU} uses it without a printed derivation.
    % Source lines: paper_v2.tex 1316--1319 and 1624--1630.
\end{lemma}

\begin{proof}\leanok
    In product coordinates the exchange matrix has entries
    $\mathbb S_{(i,j),(k,l)}=\delta_{i,l}\delta_{j,k}$, so both sides of~(\ref{eq:matrix_swap_kron_left}) have entry
    \begin{align}
        \left[\mathbb S(A\otimes B)\right]_{(i_1,i_2),(j_1,j_2)}
         & =A_{i_2,j_1}B_{i_1,j_2}
        =\left[(B\otimes A)\mathbb S\right]_{(i_1,i_2),(j_1,j_2)}.
    \end{align}
    Swapping factors $A$ and $B$ gives~(\ref{eq:matrix_swap_kron_right}).
\end{proof}

\begin{lemma}[Exchange contraction of a scaled unitary Kronecker product]
    \label{thm:matrix_trace_swap_scaled_kronecker}
    \lean{Matrix.trace_swapMatrix_mul_of_eq_smul_kronecker_conjTranspose}
    \leanok
    Let $x\in\mathbb C^{n\times n}$ be unitary, let $\sigma\in\mathbb C$, and let
    $M=\sigma\,x\otimes x^\dagger$. Then
    \begin{align}
        \tr\!\left[\mathbb SM\right] & =\sigma n.
    \end{align}
    This is the exchange-trace evaluation in
    \cite[lines~1608--1610]{Cirac2017MPU}; there $n=d^2$.
    % Source lines: paper_v2.tex 1608--1610.
\end{lemma}

\begin{proof}\leanok
    \uses{lem:matrix_trace_swap_kronecker}
    The trace is linear, so
    $\tr[\mathbb SM]=\sigma\tr[\mathbb S(x\otimes x^\dagger)]$, and
    Lemma~\ref{lem:matrix_trace_swap_kronecker} evaluates the remaining
    contraction as $n$.
\end{proof}

\begin{lemma}[Trace formula for the time-reversal sign]
    \label{eq:sigma-from-trace}
    \notready
    \discussion{5714}
    \uses{PropChiral, lem:matrix_trace_swap_kronecker}
    For a time-reversal-invariant MPU in standard form, the sign $\sigma$ of
    Lemma~\ref{PropChiral} is
    \begin{align}
        \sigma & =\frac1{d^2}\tr\!\left[\mathbb S_{12,34}
            v_{23}u_{12}u_{34}v_{23}\right],
    \end{align}
    where $\mathbb S_{12,34}$ exchanges the pairs of sites $12$ and $34$.
    % Source lines: paper_v2.tex 1588--1611.
\end{lemma}

\begin{lemma}[Admissible trace formula for the time-reversal sign]
    \label{lem:mpu_admissible_sigma_trace}
    \notready
    \discussion{5714}
    % **Scope restriction (two-site blocked comparison):** the trace formula uses
    % the blocked source data required by Lemma~\ref{prop:mpu_admissible_local_time_reversal}.
    % Documented in docs/paper-gaps/mpu_canonical_form_full_support.tex.
    \uses{prop:mpu_admissible_local_time_reversal, def:mpu_reduced_full_support_source_datum}
    For a time-reversal-invariant MPU in the standard form determined by its
    given source datum, suppose its
    one-site tensor $\mathcal U$, its two-site block $\mathcal U_2$, and
    the physical-adjoint comparison tensor $\mathcal U_2^\sharp$ are simple and
    each carries its own reduced full-support source datum. Then the sign
    $\sigma$ in
    Lemma~\ref{prop:mpu_admissible_local_time_reversal} is
    \begin{align}
        \sigma & =\frac1{d^2}\tr\!\left[\mathbb S_{12,34}v_{23}u_{12}u_{34}v_{23}\right],
    \end{align}
    where $\mathbb S_{12,34}$ exchanges the pairs of sites $12$ and $34$.
    % Source lines: paper_v2.tex 1588--1611.
\end{lemma}

\begin{lemma}[Gauge invariance of the time-reversal sign]
    \label{lem:mpu_sigma_gauge_invariant}
    \lean{MPOTensor.timeReversalTrace_gauge}
    \leanok
    Let $u,v\in\operatorname{End}(\C^d\otimes\C^d)$ and let
    $x,y\in U(d)$. Set
    \begin{align}
        u' & =(x\otimes y)u,                 &
        v' & =v(y^\dagger\otimes x^\dagger).
    \end{align}
    The four-site trace is unchanged:
    \begin{align}
        \tr\!\left[\mathbb S_{12,34}v'_{23}u'_{12}u'_{34}v'_{23}\right]
         & =\tr\!\left[\mathbb S_{12,34}v_{23}u_{12}u_{34}v_{23}\right].
    \end{align}
    In particular, the trace expression in~(\ref{eq:sigma-from-trace}) is
    invariant under the standard-form gauge.
    % Source proof: paper_v2.tex 1624--1630.
\end{lemma}

\begin{proof}\leanok
    \uses{thm:matrix_swap_kronecker_intertwine}
    With $x_i$ and $y_i$ denoting the unitaries acting on site $i$, the gated factors are
    $u'_{12}=x_1y_2u_{12}$, $u'_{34}=x_3y_4u_{34}$, and
    $v'_{23}=v_{23}y_2^\dagger x_3^\dagger$.
    Cancellation of $y_2^\dagger y_2$ and $x_3^\dagger x_3$ gives
    $v'_{23}u'_{12}u'_{34}v'_{23}
        =x_1y_4v_{23}u_{12}u_{34}v_{23}y_2^\dagger x_3^\dagger$,
    since the outer-leg operators commute with $v_{23}$. The identity
    $\mathbb S_{12,34}x_1y_4=x_3y_2\mathbb S_{12,34}$ then gives
    \begin{align}
        \tr\!\left[\mathbb S_{12,34}v'_{23}u'_{12}u'_{34}v'_{23}\right]
         & =\tr\!\left[x_3y_2\mathbb S_{12,34}
        v_{23}u_{12}u_{34}v_{23}y_2^\dagger x_3^\dagger\right]           \\
         & =\tr\!\left[\mathbb S_{12,34}v_{23}u_{12}u_{34}v_{23}\right].
    \end{align}
    The last equality follows from cyclicity of the trace and unitarity of
    $x$ and $y$.
\end{proof}

\begin{lemma}[Admissible gauge invariance of the time-reversal sign]
    \label{lem:mpu_admissible_sigma_gauge_invariant}
    \notready
    \discussion{5714}
    \uses{lem:mpu_admissible_sigma_trace}
    % Source proof: paper_v2.tex 1624--1630.
    Under the hypotheses of Lemma~\ref{lem:mpu_admissible_sigma_trace}, let $x$ and $y$ be
    unitaries and change the standard-form gauge by
    \begin{align}
        u' & =(x\otimes y)u,                 &
        v' & =v(y^\dagger\otimes x^\dagger).
    \end{align}
    Then the trace formula has the same value for $(u',v')$ as for $(u,v)$.
    Consequently the sign $\sigma$ is independent of the standard-form gauge.
\end{lemma}

\begin{proof}
    \uses{lem:mpu_admissible_sigma_trace, thm:matrix_swap_kronecker_intertwine}
    With $x_i$ and $y_i$ denoting the unitaries acting on site $i$, the gated factors are
    $u'_{12}=x_1y_2u_{12}$, $u'_{34}=x_3y_4u_{34}$, and
    $v'_{23}=v_{23}y_2^\dagger x_3^\dagger$.
    Substitution into Lemma~\ref{lem:mpu_admissible_sigma_trace}, cancellation on the
    middle legs, and
    $\mathbb S_{12,34}x_1y_4=x_3y_2\mathbb S_{12,34}$ give
    \begin{align}
        \tr\!\left[\mathbb S_{12,34}v'_{23}u'_{12}u'_{34}v'_{23}\right]
         & =\tr\!\left[x_3y_2\mathbb S_{12,34}
        v_{23}u_{12}u_{34}v_{23}y_2^\dagger x_3^\dagger\right]           \\
         & =\tr\!\left[\mathbb S_{12,34}v_{23}u_{12}u_{34}v_{23}\right].
    \end{align}
    The last equality follows from cyclicity of the trace and unitarity of
    $x$ and $y$.
\end{proof}

\begin{lemma}[Well-definedness of the time-reversal sign]
    \label{prop:sigma-well-defined}
    \notready
    \discussion{5714}
    \uses{eq:sigma-from-trace, lem:mpu_sigma_gauge_invariant}
    For a time-reversal-invariant MPU in standard form, the sign
    $\sigma=\pm1$ is independent of the gauge and is constant under continuous
    deformations that preserve time-reversal symmetry.
    % Source lines: paper_v2.tex 1613--1631.
    % The intrinsic auxiliary results below need pointwise supplied standard
    % forms and the local certificate; no continuous gate choices are needed.
    % Deducing that certificate from global symmetry remains separate.
\end{lemma}

\begin{lemma}[Extraction of the sign from a supplied local equation]
    \label{lem:mpu_supplied_time_reversal_sign}
    \lean{MPOTensor.timeReversalSign, MPOTensor.TwoSiteStandardFormData.timeReversalSign_eq_of_local_relation}
    \leanok
    \uses{def:mpu_supplied_two_site_standard_form}
    Let $d>0$ and let $W$ have a supplied two-site standard form with square
    gates $u,v$ on $\mathbb C^d\otimes\mathbb C^d$. Define
    $s(u,v)=d^{-2}\tr[\mathbb S_{12,34}v_{23}(u_{12}\otimes u_{34})v_{23}]$.
    If a unitary $x$ and a scalar $\sigma$ satisfy
    \begin{align}
        v_{23}(u_{12}\otimes u_{34})
         & =\sigma(x\otimes x^\dagger)v_{23}^\dagger,
    \end{align}
    then $s(u,v)=\sigma$. This is the local-equation extraction in
    \cite[lines~1595--1610]{Cirac2017MPU} for the supplied standard form.
\end{lemma}

\begin{proof}\leanok
    \uses{thm:matrix_trace_swap_scaled_kronecker}
    Multiply the local equation by $v_{23}$ and cancel
    $v_{23}^\dagger v_{23}=\Id$. Exchange of the two physical pairs gives
    $\tr[\mathbb S(x\otimes x^\dagger)]=\tr[xx^\dagger]=d^2$.
    Division by the nonzero scalar $d^2$ gives the result.
\end{proof}

\begin{lemma}[Intrinsic trace and independence of standard-form representatives]
    \label{lem:mpu_source_time_reversal_sign_independent}
    \lean{MPOTensor.timeReversalPeriodicTrace, MPOTensor.intrinsicTimeReversalSign, MPOTensor.TwoSiteStandardFormData.timeReversalTrace_eq_periodicTrace, MPOTensor.TwoSiteStandardFormData.timeReversalSign_eq_intrinsic, MPOTensor.TwoSiteStandardFormData.timeReversalSign_eq_of_mpo_two_eq}
    \leanok
    \uses{def:mpu_supplied_two_site_standard_form}
    Let $W$ have a supplied two-site standard form with square gates $u,v$
    on $\mathbb C^d\otimes\mathbb C^d$. For exchange $P$ of the two physical
    blocks, define $s(W)=d^{-2}\tr[P U_W^{(2)}]$. Then $s(u,v)=s(W)$.
    Consequently two supplied standard forms whose two-block
    periodic operators agree have the same normalized sign expression,
    including representatives in different bond dimensions. No canonical
    source-factor comparison is required. This gives the independence
    assertion of \cite[lines~1613--1627]{Cirac2017MPU} for the supplied
    square standard forms.
\end{lemma}

\begin{proof}\leanok
    \uses{thm:mpu_supplied_two_site_periodic_circuit}
    The shifted layer on two blocks is $P v_{23}P v_{23}$, so
    $U_W^{(2)}=P v_{23}P v_{23}(u\otimes u)$. Since $P^2=\Id$,
    cyclicity gives
    \begin{align}
        \tr[P U_W^{(2)}]
         & =\tr[v_{23}P v_{23}(u\otimes u)]  \\
         & =\tr[P v_{23}(u\otimes u)v_{23}].
    \end{align}
    Division by $d^2$ gives the identity and independence.
\end{proof}

\begin{lemma}[Constancy along tensor paths with pointwise local certificates]
    \label{lem:mpu_supplied_time_reversal_sign_constant}
    \lean{MPOTensor.continuous_intrinsicTimeReversalSign, MPOTensor.intrinsicTimeReversalSign_eq_on_path}
    \leanok
    \uses{lem:mpu_supplied_time_reversal_sign, def:mpu_supplied_two_site_standard_form}
    Let $W(t)$ be a continuous path of tensors for $t\in[0,1]$.
    At each $t$, suppose there is a supplied square two-site standard form
    and a local time-reversal certificate as in
    Lemma~\ref{lem:mpu_supplied_time_reversal_sign}, with
    $\sigma(t)\in\{1,-1\}$. These choices are pointwise; no continuity
    of the gates, source factors, unitary or sign witnesses is required.
    Then $s(W(t))=s(W(0))$ throughout the interval.
    This is the intrinsic-trace argument for
    \cite[lines~1628--1630]{Cirac2017MPU}. Deducing the pointwise local
    certificate from global time reversal remains the separate local
    characterization.
\end{lemma}

\begin{proof}\leanok
    \uses{lem:mpu_supplied_time_reversal_sign, lem:mpu_source_time_reversal_sign_independent}
    The finite-chain operator is polynomial in the tensor entries,
    so $s(W(t))$ is continuous. At each point the local equation and
    intrinsic identity give $s(W(t))\in\{1,-1\}$. A continuous function
    from the connected interval to this two-point set is constant.
\end{proof}

\begin{lemma}[Admissible well-definedness of the time-reversal sign]
    \label{prop:mpu_admissible_sigma_well_defined}
    \notready
    \discussion{5714}
    % **Scope restriction (admissible deformation):** the deformation carries the
    % path datum of Definition~\ref{def:mpu_admissible_equivalence_datum}.
    % Documented in docs/paper-gaps/mpu_canonical_form_full_support.tex.
    \uses{lem:mpu_admissible_sigma_gauge_invariant, lem:mpu_admissible_symmetry_path_criterion, def:mpu_reduced_full_support_source_datum, def:mpu_admissible_equivalence_datum}
    Suppose a time-reversal-preserving path carries an admissible equivalence
    datum of blocking length one. At every path point, require the one-site
    tensor, its two-site block, and the physical-adjoint comparison of
    that two-site block to be simple and to carry continuously varying reduced
    full-support source data. Then the sign
    $\sigma=\pm1$ is independent of the admissible standard-form gauge and is
    constant along the path.
    % Source lines: paper_v2.tex 1613--1631.
\end{lemma}

\begin{lemma}[Scalar parity law for the blocked sign]
    \label{thm:mpu_blocked_sign_parity_scalar}
    \lean{TNLean.Algebra.sq_mul_eq_one_of_sign,
        TNLean.Algebra.blockedSign_eq_first_of_odd,
        TNLean.Algebra.blockedSign_eq_second_of_even}
    \leanok
    Let $\sigma^{(1)},\sigma^{(2)}\in\{1,-1\}$ and put
    $\zeta=\sigma^{(1)}\sigma^{(2)}$. For every $k\geq1$,
    \begin{align}
        \sigma^{(1)}\zeta^{k-1}
         & =\begin{cases}
                \sigma^{(1)}, & k\text{ odd},  \\
                \sigma^{(2)}, & k\text{ even}.
            \end{cases}
    \end{align}
    In \cite[lines~1687 and 1730--1732]{Cirac2017MPU}, this alternating formula is the scalar
    conclusion.
    % Source lines: paper_v2.tex 1687 and 1730--1732.
\end{lemma}

\begin{proof}\leanok
    Both signs square to one, hence $\zeta^2=1$. For odd $k$ the exponent $k-1$
    is even, so $\zeta^{k-1}=1$ and the product is $\sigma^{(1)}$. For even
    $k\geq2$ the exponent $k-1$ is odd, so $\zeta^{k-1}=\zeta$ and the product
    is $(\sigma^{(1)})^2\sigma^{(2)}=\sigma^{(2)}$.
\end{proof}

\begin{lemma}[Blocking parity of the time-reversal sign]
    \label{prop:sigma1-and-sigma2-equal}
    \notready
    \discussion{5714}
    \uses{prop:sigma-well-defined, PropChiral, lem:mpu_blocking}
    Let $\mathcal U$ be a time-reversal-invariant MPU in standard form, and let
    $\sigma^{(k)}$ be the sign obtained from a standard form of its $k$-site
    block. Then
    \begin{align}
        \sigma^{(k)}
         & =\begin{cases}
                \sigma^{(1)}, & k\text{ odd},  \\
                \sigma^{(2)}, & k\text{ even}.
            \end{cases}
    \end{align}
    % Source lines: paper_v2.tex 1633--1733.
\end{lemma}

\begin{proof}
    \uses{lem:mpu_sigma_gauge_invariant, PropChiral, lem:mpu_blocking,
        thm:mpu_blocked_sign_parity_scalar}
    Gauge invariance permits the standard blocked representatives
    of the source proof, namely
    \begin{align}
        u^{(2)} & =v_{23}u_{12}u_{34}, & v^{(2)} & =v_{45},
    \end{align}
    with $\zeta=\sigma^{(2)}/\sigma^{(1)}$. The one- and two-site local
    time-reversal relations give the unwinding identity
    \begin{align}
        \zeta\big[(\Id\otimes v\otimes\Id)(u\otimes x)\big]^\dagger
         & =(\Id\otimes v\otimes\Id)(x^\dagger\otimes u).
    \end{align}
    Applying this identity successively through a block of length $k$ gives
    \begin{align}
        \sigma^{(k)}
         & =\sigma^{(1)}\zeta^{k-1}.
    \end{align}
    Since each sign is $\pm1$, this equals $\sigma^{(1)}$ for odd $k$ and
    $\sigma^{(2)}$ for even $k$.
    % Source lines: paper_v2.tex 1657--1732.
\end{proof}

\begin{lemma}[Admissible blocking parity of the time-reversal sign]
    \label{prop:mpu_admissible_sigma_blocking_parity}
    \notready
    \discussion{5714}
    % **Scope restriction (blocked source data):** every blocked tensor used to
    % define $\sigma^{(k)}$, including its two-site comparison tensor, carries
    % reduced full-support source data. No preservation under blocking is
    % asserted. Documented in
    % docs/paper-gaps/mpu_canonical_form_full_support.tex.
    \uses{lem:mpu_admissible_sigma_gauge_invariant, prop:mpu_admissible_local_time_reversal, lem:mpu_blocking, def:mpu_reduced_full_support_source_datum}
    For every blocking length $k$ under consideration, suppose the endpoint
    tensor $\mathcal U_k$, its two-site block $\mathcal U_{2k}$, and the
    physical-adjoint comparison tensor $\mathcal U_{2k}^\sharp$ are simple and
    each carries its own reduced full-support source datum. Let $\sigma^{(k)}$
    be the sign of the resulting standard form. Then
    \begin{align}
        \sigma^{(k)}
         & =\begin{cases}
                \sigma^{(1)}, & k\text{ odd},  \\
                \sigma^{(2)}, & k\text{ even}.
            \end{cases}
    \end{align}
    % Source lines: paper_v2.tex 1633--1733.
\end{lemma}

\begin{proof}
    \uses{lem:mpu_admissible_sigma_gauge_invariant, prop:mpu_admissible_local_time_reversal, lem:mpu_blocking, thm:mpu_blocked_sign_parity_scalar}
    Gauge invariance permits the standard blocked representatives
    of the source proof, namely
    \begin{align}
        u^{(2)} & =v_{23}u_{12}u_{34}, & v^{(2)} & =v_{45},
    \end{align}
    with $\zeta=\sigma^{(2)}/\sigma^{(1)}$. The one- and two-site local
    time-reversal relations give the unwinding identity
    \begin{align}
        \zeta\big[(\Id\otimes v\otimes\Id)(u\otimes x)\big]^\dagger
         & =(\Id\otimes v\otimes\Id)(x^\dagger\otimes u).
    \end{align}
    Applying this identity successively through a block of length $k$ gives
    \begin{align}
        \sigma^{(k)}
         & =\sigma^{(1)}\zeta^{k-1}.
    \end{align}
    Since each sign is $\pm1$, this equals $\sigma^{(1)}$ for odd $k$ and
    $\sigma^{(2)}$ for even $k$.
    % Source lines: paper_v2.tex 1657--1732.
\end{proof}

\begin{corollary}[Necessary time-reversal invariants]
    \label{cor:sigma-equivalence-necessary}
    \notready
    \discussion{5714}
    \uses{prop:sigma1-and-sigma2-equal, def:strictly-equivalent-symmetry, def:equivalent-symmetry}
    Strictly equivalent time-reversal-invariant MPUs have the same
    $\sigma^{(1)}$ and $\sigma^{(2)}$. Equivalent time-reversal-invariant MPUs
    have the same $\sigma^{(2)}$.
    % Source lines: paper_v2.tex 1735--1743.
\end{corollary}

\begin{corollary}[Admissible time-reversal invariants of equivalent paths]
    \label{cor:mpu_admissible_sigma_equivalence_necessary}
    \notready
    \discussion{5714}
    % **Scope restriction (admissible symmetry equivalence):** the endpoint
    % blockings and symmetry paths carry the data of
    % Definition~\ref{def:mpu_admissible_equivalence_datum}. The non-strict
    % conclusion concerns only the supplied ancilla-enlarged blocked endpoints.
    % Documented in docs/paper-gaps/mpu_canonical_form_full_support.tex.
    \uses{prop:mpu_admissible_sigma_blocking_parity, def:mpu_admissible_strict_symmetry_equivalence, def:mpu_admissible_symmetry_equivalence, def:mpu_reduced_full_support_source_datum, def:mpu_admissible_equivalence_datum}
    Let two time-reversal-invariant MPUs be compared by an admissible symmetry
    path $p\mapsto\mathcal W(p)$ whose path datum has blocking length one.
    Suppose that the five path families $\mathcal W(p)$,
    $\mathcal W(p)_2$, $\mathcal W(p)_4$,
    $\mathcal W(p)_2^\sharp$, and $\mathcal W(p)_4^\sharp$ are simple and carry
    continuously varying reduced full-support source data. Then admissible
    strict equivalence requires equality of both $\sigma^{(1)}$ and
    $\sigma^{(2)}$.

    For admissible equivalence after adjoining ancillas and applying the source
    blocking, let $p\mapsto\mathcal W_{\mathrm{anc}}(p)$ be the
    ancilla-enlarged blocked connecting path itself, not a reduced or equivalent
    representative, and regard it as a length-one admissible path. Require
    continuous reduced full-support source data along
    $\mathcal W_{\mathrm{anc}}(p)$, its two- and four-site blocks, and
    the physical-adjoint comparisons of those two blocks. The conclusion is
    only
    \begin{align}
        \sigma^{(2)}(\mathcal W_{\mathrm{anc}}(0))
         & =\sigma^{(2)}(\mathcal W_{\mathrm{anc}}(1)).
    \end{align}
    No equality between $\sigma^{(2)}$ of the original two MPUs is asserted.
    Such a conclusion requires preservation under adjoining identity ancillas,
    which is not established here.
    % Tracked by issue #6138.
    % Source lines: paper_v2.tex 1735--1743.
\end{corollary}

\begin{lemma}[Conjugation class and the MPS sign]
    \label{prop:mps-class-eq-mpu-class-for-conjugation}
    \notready
    \discussion{5714}
    \uses{corconj, cor:equiv-under-conjugation}
    For a conjugation-invariant MPU, the MPS sign
    $G\overline G=\pm\Id$ agrees with Case I and Case II, respectively, of
    Lemma~\ref{corconj}.
    % Source lines: paper_v2.tex 1791--1805.
\end{lemma}

\begin{lemma}[Admissible conjugation class and the MPS sign]
    \label{prop:mpu_admissible_mps_conjugation_class}
    \notready
    \discussion{5714}
    % **Scope restriction (admissible conjugation class):** the standard forms,
    % blockings, and comparison paths carry the required source data. Documented
    % in docs/paper-gaps/mpu_canonical_form_full_support.tex.
    \uses{prop:mpu_admissible_conjugation_standard_forms, cor:mpu_admissible_equiv_under_conjugation, def:mpu_reduced_full_support_source_datum, def:mpu_admissible_equivalence_datum}
    Let $\mathcal U$ be conjugation invariant. Suppose $\mathcal U$ and its
    conjugate comparison tensor are simple and carry their own reduced
    full-support source data. Let $p\mapsto\mathcal W(p)$ be the explicit
    standard-form interpolation from $\mathcal U$ to its Case I or Case II
    representative. Suppose this interpolation carries an admissible
    conjugation-symmetry path datum of blocking length one: the path
    $\mathcal W(p)$ and its fixed conjugate, physical-adjoint, transposed, and
    conjugation-image comparison paths are simple and carry continuously varying
    source data. Then the sign
    $G\overline G=\pm\Id$ in the MPS classification agrees with Case I and
    Case II, respectively, of
    Lemma~\ref{prop:mpu_admissible_conjugation_standard_forms}.
    % Source lines: paper_v2.tex 1791--1805.
\end{lemma}

\begin{lemma}[Time-reversal class and the MPS sign]
    \label{prop:dagger-mps-class-and-mpu-class-equal}
    \notready
    \discussion{5714}
    \uses{PropChiral, prop:sigma1-and-sigma2-equal}
    For a time-reversal-invariant MPU in standard form, the MPS sign
    $G\overline G=\pm\Id$ equals the sign $\sigma^{(2)}$ obtained after
    blocking two sites.
    % Source lines: paper_v2.tex 1807--1829.
\end{lemma}

\begin{lemma}[Admissible time-reversal class and the MPS sign]
    \label{prop:mpu_admissible_mps_time_reversal_class}
    \notready
    \discussion{5714}
    % **Scope restriction (one-, two-, and four-site source data):** the tensors
    % used to define both $\sigma^{(1)}$ and $\sigma^{(2)}$ carry their own data.
    % Documented in docs/paper-gaps/mpu_canonical_form_full_support.tex.
    \uses{prop:mpu_admissible_local_time_reversal, prop:mpu_admissible_sigma_blocking_parity, def:mpu_reduced_full_support_source_datum, def:mpu_admissible_equivalence_datum}
    Let an MPU in the standard form determined by its given source datum
    be invariant under time reversal.
    Suppose $\mathcal U$, $\mathcal U_2$, and $\mathcal U_4$, together with the
    physical-adjoint comparison tensors $\mathcal U_2^\sharp$ and
    $\mathcal U_4^\sharp$, are simple and each carries its own reduced
    full-support source datum. Then the MPS sign $G\overline G=\pm\Id$ equals
    the sign $\sigma^{(2)}$.
    % Source lines: paper_v2.tex 1807--1829.
\end{lemma}
